\chapter{The Stability of Superposed Fluids: The Rayleigh--Taylor Instability}
\label{ch:10}

\section{Introduction}\label{sec:90}

In the last three chapters we considered problems in which instabilities arose, principally, from adverse distributions of angular momentum. In the present chapter and in the next, we shall be concerned with instabilities which arise from different causes: they are of two kinds. The \emph{first} derives from the character of the equilibrium of an incompressible heavy fluid of variable density (i.e.\ of a heterogeneous fluid). An important special case in this connexion is that of two fluids of different densities superposed one over the other (or accelerated towards each other); the instability of the plane interface between the two fluids, when it occurs (particularly in the second context), is called \emph{Rayleigh--Taylor} instability. The second type of instability arises when the different layers of a stratified heterogeneous fluid are in relative horizontal motion. An important special case in this latter connexion is when two superposed fluids flow one over the other with a relative horizontal velocity; the instability of the plane interface between the two fluids, when it occurs in this instance, is called \emph{Kelvin--Helmholtz instability}.

In this chapter we shall be concerned with Rayleigh--Taylor instability.

\section{The character of the equilibrium of a stratified heterogeneous fluid. The perturbation equations}\label{sec:91}

A static state, in which an incompressible fluid of variable density is arranged in horizontal strata and the pressure $p$ and the density $\rho$ are functions of the vertical coordinate $z$ only, is clearly a kinematically realizable one. The character of the equilibrium of this initial static state can be determined, as usual, by supposing that the system is slightly disturbed and then by following its further evolution.

Let the actual density at any point $(x,y,z)$ as a result of the disturbance be $\rho+\delta\rho$; and let $\delta p$ denote the corresponding increment in the pressure. And finally, let $u_i$ ($i=1,2,3$) denote the components of the velocity (considered small). The relevant equations of motion and continuity, then, are
\begin{equation}
\label{eq:10-1}
\rho\frac{\partial u_i}{\partial t} = -\frac{\partial}{\partial x_i}\,\delta p + \frac{\partial}{\partial x_j}\,p_{ij} - g\,\delta\rho\,\lambda_i
\end{equation}
and
\begin{equation}
\label{eq:10-2}
\frac{\partial u_i}{\partial x_i} = 0,
\end{equation}
where $g$ is the acceleration due to gravity and $\boldsymbol{\lambda}$ is a unit vector in the direction of the vertical. In equation~\eqref{eq:10-1} $p_{ij}$ is the viscous stress-tensor given by
\begin{equation}
\label{eq:10-3}
p_{ij} = \mu\!\left(\frac{\partial u_i}{\partial x_j}+\frac{\partial u_j}{\partial x_i}\right)\!,
\end{equation}
where $\mu$ (which can be a function of $z$) denotes the coefficient of viscosity. In addition to equations~\eqref{eq:10-1} and~\eqref{eq:10-2}, we have the equation,
\begin{equation}
\label{eq:10-4}
\frac{\partial}{\partial t}\,\delta\rho + u_i\frac{\partial\rho}{\partial x_i} = 0,
\end{equation}
which ensures that the density of every particle remains unchanged as we follow it with its motion.

Inserting equation~\eqref{eq:10-3} in equation~\eqref{eq:10-1} and making use of the solenoidal character of $u_i$, we find that the equation of motion reduces to the form
\begin{equation}
\label{eq:10-5}
\rho\frac{\partial u_i}{\partial t} = -\frac{\partial}{\partial x_i}\,\delta p + \mu\nabla^2 u_i + \left(\frac{\partial u_i}{\partial x_j}+\frac{\partial u_j}{\partial x_i}\right)\frac{\partial\mu}{\partial x_j} - g\,\delta\rho\,\lambda_i
\end{equation}
or, since
\begin{equation}
\label{eq:10-6}
\frac{\partial\mu}{\partial x_j} = \lambda_j\frac{d\mu}{dz},
\end{equation}
we have
\begin{equation}
\label{eq:10-7}
\rho\frac{\partial u_i}{\partial t} = -\frac{\partial}{\partial x_i}\,\delta p + \mu\nabla^2 u_i + \left(\frac{\partial w}{\partial x_i}+\frac{\partial u_i}{\partial z}\right)\frac{d\mu}{dz} - g\,\delta\rho\,\lambda_i,
\end{equation}
where $w$ $(= \lambda_i u_i)$ is the $z$-component of the velocity.

It is now convenient to have the equations for the components of the velocity $u$, $v$, and $w$ (in the $x$-, $y$-, and $z$-directions, respectively) written out explicitly. We have
\begin{equation}
\label{eq:10-8}
\rho\frac{\partial u}{\partial t} = -\frac{\partial}{\partial x}\,\delta p + \mu\nabla^2 u + \left(\frac{\partial w}{\partial x}+\frac{\partial u}{\partial z}\right)\frac{d\mu}{dz},
\end{equation}
\begin{equation}
\label{eq:10-9}
\rho\frac{\partial v}{\partial t} = -\frac{\partial}{\partial y}\,\delta p + \mu\nabla^2 v + \left(\frac{\partial w}{\partial y}+\frac{\partial v}{\partial z}\right)\frac{d\mu}{dz},
\end{equation}
\begin{equation}
\label{eq:10-10}
\rho\frac{\partial w}{\partial t} = -\frac{\partial}{\partial z}\,\delta p + \mu\nabla^2 w + 2\frac{\partial w}{\partial z}\frac{d\mu}{dz} - g\,\delta\rho,
\end{equation}
\begin{equation}
\label{eq:10-11}
\frac{\partial u}{\partial x}+\frac{\partial v}{\partial y}+\frac{\partial w}{\partial z} = 0,
\end{equation}
and
\begin{equation}
\label{eq:10-12}
\frac{\partial}{\partial t}\,\delta\rho = -w\frac{d\rho}{dz}.
\end{equation}

Analysing the disturbance into normal modes, we seek solutions whose dependence on $x$, $y$, and $t$ is given by
\begin{equation}
\label{eq:10-13}
\exp(ik_x\,x + ik_y\,y + nt),
\end{equation}
where $k_x$, $k_y$, and $n$ are constants. For solutions having this dependence on $x$, $y$, and $t$, equations~\eqref{eq:10-8}--\eqref{eq:10-12} become
\begin{equation}
\label{eq:10-14}
ik_x\,\delta p = -n\rho u + \mu(D^2-k^2)u + (D\mu)(ik_x\,w+Du),
\end{equation}
\begin{equation}
\label{eq:10-15}
ik_y\,\delta p = -n\rho v + \mu(D^2-k^2)v + (D\mu)(ik_y\,w+Dv),
\end{equation}
\begin{equation}
\label{eq:10-16}
D\,\delta p = -n\rho w + \mu(D^2-k^2)w + 2(D\mu)(Dw) - g\,\delta\rho,
\end{equation}
\begin{equation}
\label{eq:10-17}
ik_x\,u + ik_y\,v = -Dw,
\end{equation}
and
\begin{equation}
\label{eq:10-18}
n\,\delta\rho = -wDp,
\end{equation}
where
\begin{equation}
\label{eq:10-19}
k^2 = k_x^2+k_y^2 \quad\text{and}\quad D = d/dz.
\end{equation}
Multiplying equations~\eqref{eq:10-14} and~\eqref{eq:10-15} by $-ik_x$ and $-ik_y$, respectively, adding, and making use of equation~\eqref{eq:10-17}, we obtain
\begin{equation}
\label{eq:10-20}
k^2\delta p = [-n\rho+\mu(D^2-k^2)]Dw + (D\mu)(D^2+k^2)w;
\end{equation}
whereas by combining equations~\eqref{eq:10-16} and~\eqref{eq:10-18} we have
\begin{equation}
\label{eq:10-21}
D\,\delta p = -n\rho w + \mu(D^2-k^2)w + 2(D\mu)(Dw) + \frac{g}{n}(D\rho)w.
\end{equation}

Now eliminating $\delta p$ between equations~\eqref{eq:10-20} and~\eqref{eq:10-21}, we obtain
\begin{align}
\label{eq:10-22}
D\!\left[\rho-\frac{\mu}{n}(D^2-k^2)\right]\!Dw - \frac{1}{n}(D\mu)(D^2+k^2)w \nonumber\\
= k^2\!\left[\rho-\frac{\mu}{n}(D^2-k^2)\right]\!w - \frac{g}{n^2}(D\rho)w. \qquad\quad
\end{align}

This is the required equation. If we suppose that the fluid is confined between two rigid planes, then the boundary conditions are
\begin{equation}
\label{eq:10-23}
w = 0 \quad\text{and}\quad Dw = 0 \quad\text{on the bounding surfaces,}
\end{equation}
where the latter condition ensures the vanishing of the horizontal components of the velocity.

\subsection*{(a) \textit{Allowance for surface tension at interfaces between fluids}}

In deriving equation~\eqref{eq:10-22} we have argued as though the density were continuously variable. The case when fluids of different densities are superposed one over the other can be formally included in the analysis by admitting discontinuities in the distribution of the density. This, however, is not entirely sufficient: for the interfaces between fluids will be subject to forces arising from surface tension; and we have not allowed for this possibility in deriving equation~\eqref{eq:10-22}. We shall now consider the modifications which are necessary to allow for such forces coming into play at places where the density changes discontinuously.

Now the effect of surface tension is that, for equilibrium, the normal stresses acting on the two sides ($+$ and $-$) of an element of surface $dS$ separating two fluids must differ by an amount depending on the \emph{surface tension} $T$ and given by
\begin{equation}
\label{eq:10-24}
[(-p\delta_{ij}+p_{ij})-(-p\delta_{ij}+p_{ij})]N_j = -T\!\left(\frac{1}{R_1}+\frac{1}{R_2}\right)\!,
\end{equation}
where $N_i$ is the unit outward normal to $dS$, and $R_1$ and $R_2$ are the principal radii of curvature of $dS$. The convention regarding signs is that $R_1$ and $R_2$ are to be considered positive if the corresponding centres of curvature lie on the side of $dS$ to which the plus sign refers.

Returning to the problem we were considering, suppose that discontinuities in $\rho$ occur at some preassigned levels, $z_s$ ($s = 1, z_2, z_3, \ldots$). In the perturbed state these plane surfaces will become slightly deformed. Let the surfaces of separation be then defined by
\begin{equation}
\label{eq:10-25}
z_s+\delta z_s(x,y,t).
\end{equation}
On account of~\eqref{eq:10-24}, the discontinuity in the normal stresses required for equilibrium is
\begin{equation}
\label{eq:10-26}
T\!\left(\frac{\partial^2}{\partial x^2}+\frac{\partial^2}{\partial y^2}\right)\!\delta z_s.
\end{equation}
Therefore, to allow for the effects of surface tension, we must include on the right-hand side of equation~\eqref{eq:10-10} by the additional term
\begin{equation}
\label{eq:10-27}
\sum_s\left[T_s\!\left(\frac{\partial^2}{\partial x^2}+\frac{\partial^2}{\partial y^2}\right)\!\delta z_s\right]\delta(z-z_s),
\end{equation}
where $\delta(z-z_s)$ denotes Dirac's $\delta$-function. Thus, we replace equation~\eqref{eq:10-10} by
\begin{align}
\label{eq:10-28}
\rho\frac{\partial w}{\partial t} = -\frac{\partial}{\partial z}\,\delta p + \mu\nabla^2 w + 2\frac{\partial w}{\partial z}\frac{d\mu}{dz} - g\,\delta\rho + {} \nonumber\\
{} + \sum_s T_s\!\left(\frac{\partial^2}{\partial x^2}+\frac{\partial^2}{\partial y^2}\right)\!\delta z_s\,\delta(z-z_s).
\end{align}
Equations~\eqref{eq:10-8} and~\eqref{eq:10-9} remain unaffected.

Equation~\eqref{eq:10-21}, appropriate for solutions having an $(x,y,t)$-dependence given by~\eqref{eq:10-13}, is now replaced by
\begin{equation}
\label{eq:10-29}
D\,\delta p = -n\rho w + \mu(D^2-k^2)w + 2(D\mu)(Dw) + \frac{g}{n}(D\rho)w - k^2\sum_s (T_s\delta z_s)\delta(z-z_s).
\end{equation}
In this equation $\delta z_s$ can be expressed in terms of the normal component of the velocity $w$ at $z_s$ by observing that
\begin{equation}
\label{eq:10-30}
\frac{d}{dt}\,\delta z_s = w_s
\end{equation}
and that, therefore,
\begin{equation}
\label{eq:10-31}
\delta z_s = w_s/n.
\end{equation}
The required generalization of equation~\eqref{eq:10-21} is thus:
\begin{equation}
\label{eq:10-32}
D\,\delta p = -n\rho w + \mu(D^2-k^2)w + 2(D\mu)(Dw) + \frac{g}{n}(D\rho)w - \frac{k^2}{n}\sum_s (T_s w_s)\delta(z-z_s).
\end{equation}

To a certain extent, equation~\eqref{eq:10-32} has been obtained by a formal procedure. And the fact that the equation involves $\delta$-functions means that to interpret the equation correctly at a point of discontinuity, we must integrate the equation, across the interface, over an infinitesimal element of $z$ including the point of discontinuity. Before we effect such an integration, it is necessary to specify which quantities should be continuous and which bounded at an interface between two fluids. Clearly, all three components of the velocity, as well as the tangential viscous stresses, must be continuous. The continuity of $Dw$ follows from equation~\eqref{eq:10-17} and the continuity of $u$ and $v$. To obtain the condition which ensures the continuity of the tangential viscous stresses, we first observe that
\begin{equation}
\label{eq:10-33}
p_{xz} = \mu\!\left(\frac{\partial u}{\partial z}+\frac{\partial w}{\partial x}\right) = \mu(Du+ik_x\,w)
\end{equation}
and
\begin{equation}
\label{eq:10-34}
p_{yz} = \mu\!\left(\frac{\partial v}{\partial z}+\frac{\partial w}{\partial y}\right) = \mu(Dv+ik_y\,w).
\end{equation}
From these equations we obtain
\begin{equation}
\label{eq:10-35}
i(k_x\,p_{xz}+k_y\,p_{yz}) = \mu[D(ik_x\,u+ik_y\,v)-k^2w],
\end{equation}
or, making use of equation~\eqref{eq:10-17}, we have
\begin{equation}
i(k_x\,p_{xz}+k_y\,p_{yz}) = -\mu(D^2+k^2)w.
\end{equation}
Hence, $\mu(D^2+k^2)w$ must remain continuous across an interface. Altogether, the conditions at an interface are
\begin{equation}
\label{eq:10-36}
\textit{$w$, $Dw$, and $\mu(D^2+k^2)w$ must be continuous across an interface between two fluids.}
\end{equation}
The last of these conditions requires, in particular, that $D^2w$ \emph{must be bounded at an interface}.

In view of the continuity of $w$ and $Dw$, and of the boundedness of $D^2w$ at an interface, the result of integrating equation~\eqref{eq:10-32} over an infinitesimal element of $z$ including $z_s$ (for example) is
\begin{equation}
\label{eq:10-37}
\Delta_s(\delta p) = 2(Dw)_s\,\Delta_s(\mu)+\frac{g}{n}\,w_s\,\Delta_s(\rho)-\frac{k^2}{n}\,T_s\,w_s,
\end{equation}
where
\begin{equation}
\label{eq:10-38}
\Delta_s(f) = f(z_s+0)-f(z_s-0)
\end{equation}
is the jump which a quantity $f$ experiences at $z = z_s$, and the subscript $s$ distinguishes the value a quantity, known to be continuous at an interface, takes at $z = z_s$. On the other hand, according to equation~\eqref{eq:10-20} (the validity of which is unaffected by surface tension),
\begin{equation}
\label{eq:10-39}
k^2\Delta_s(\delta p) = \Delta_s\{[-n\rho+\mu(D^2-k^2)]Dw + [(D^2+k^2)w]D\mu\}.
\end{equation}
Combining equations~\eqref{eq:10-37} and~\eqref{eq:10-39}, we have
\begin{align}
\label{eq:10-40}
\Delta_s\!\left\{\!\left[\rho-\frac{\mu}{n}(D^2-k^2)\right]\!Dw - \frac{1}{n}[(D^2+k^2)w]\,D\mu\right\} \nonumber\\
= -\frac{k^2}{n^2}\,[g\,\Delta_s(\rho)-k^2T_s]w_s - \frac{2k^2}{n}\,(Dw)_s\,\Delta_s(\mu). \qquad
\end{align}
This equation may be considered as a boundary condition (in addition to those stated in~\eqref{eq:10-36}) which must be satisfied at an interface between two fluids.

With the interpretation of~\eqref{eq:10-37} of equation~\eqref{eq:10-32}, we can eliminate $\delta p$ between this equation and~\eqref{eq:10-20} to obtain
\begin{align}
\label{eq:10-41}
D\!\left\{\!\left[\rho-\frac{\mu}{n}(D^2-k^2)\right]\!Dw - \frac{1}{n}(D\mu)(D^2+k^2)w\right\} \nonumber\\
= k^2\!\left\{\!\left[-\frac{n}{g}(D\rho)+\frac{k^2}{n}\sum_s T_s\,\delta(z-z_s)\right]\!w + {} \right. \nonumber\\
\left. {} + \left[\rho-\frac{\mu}{n}(D^2-k^2)\right]\!w - \frac{g}{n^2}(D\rho)(Dw)\right\}.
\end{align}
This is the required generalization of equation~\eqref{eq:10-22}. The correct interpretation of equation~\eqref{eq:10-41} is that for $z \neq z_s$ ($s = 1, 2, \ldots$) the equation is identical with~\eqref{eq:10-22}, while at $z = z_s$, it is equivalent to~\eqref{eq:10-40}.

\section{The inviscid case}\label{sec:92}

We shall first consider the case when the fluid is inviscid and the stratification of density is moreover continuous. Then equation~\eqref{eq:10-22} becomes
\begin{equation}
\label{eq:10-42}
D(\rho Dw) - \rho k^2w = -\frac{k^2}{n^2}\,g(D\rho)w,
\end{equation}
and the boundary conditions require only that
\begin{equation}
\label{eq:10-43}
w = 0 \quad\text{on the boundaries.}
\end{equation}

The characteristic value problem presented by equations~\eqref{eq:10-42} and~\eqref{eq:10-43} is of the standard type to which the classical Sturmian theory applies. We conclude: the characteristic values of $n^2$ are all positive if $D\rho$ is everywhere positive; they are all negative if $D\rho$ is everywhere negative; and if $D\rho$ should change sign anywhere inside the fluid, then both positive and negative characteristic values for $n^2$ exist. Thus, \emph{the necessary and sufficient condition that a stratified heterogeneous fluid be stable is that $D\rho$ should be negative everywhere; moreover, if $D\rho$ should be positive anywhere inside the fluid, the stratification is unstable.}

It is also clear that the characteristic value problem is self-adjoint and allows a variational formulation. The basis for the latter is the integral relation,
\begin{equation}
\label{eq:10-44}
\frac{n^2}{k^2} = \frac{g\int(D\rho)w^2\,dz}{\int[\rho(Dw)^2+k^2w^2]\,dz},
\end{equation}
which one readily obtains. We observe the complete similarity of equation~\eqref{eq:10-44} with equation VII~(60); and the entirely equivalent roles of $D\rho$ and of $\theta(r)$ for the two problems is apparent.

\subsection*{(a) \textit{The case of two uniform fluids of constant density separated by a horizontal boundary}}

Consider the case of two uniform fluids of densities $\rho_1$ and $\rho_2$ separated by a horizontal boundary at $z = 0$. For both regions of the fluid, the general equation~\eqref{eq:10-42} reduces to
\begin{equation}
\label{eq:10-45}
(D^2-k^2)w = 0
\end{equation}
of which the general solution is
\begin{equation}
\label{eq:10-46}
w = Ae^{+kz}+Be^{-kz}.
\end{equation}
Since $w$ must vanish both when $z \to -\infty$ (in the lower fluid) and $z \to +\infty$ (in the upper fluid), we must suppose that
\begin{equation}
\label{eq:10-47}
w_1 = Ae^{+kz} \quad (z < 0)
\end{equation}
and
\begin{equation}
\label{eq:10-48}
w_2 = Ae^{-kz} \quad (z > 0),
\end{equation}
where we have chosen the same constant $A$ in the solutions for $z > 0$ and $z < 0$ to ensure the continuity of $w$ across the interface at $z = 0$. The remaining condition~\eqref{eq:10-40}, in the inviscid case, requires
\begin{equation}
\label{eq:10-49}
\Delta_0(\rho Dw) = -\frac{k^2}{n^2}\,[g(\rho_2-\rho_1)-k^2T]w_0.
\end{equation}
Applying the condition~\eqref{eq:10-49} to the solution~\eqref{eq:10-47} and~\eqref{eq:10-48}, we obtain
\begin{equation}
\label{eq:10-50}
-k(\rho_2+\rho_1) = -\frac{k^2}{n^2}[g(\rho_2-\rho_1)-k^2T]
\end{equation}
or
\begin{equation}
\label{eq:10-51}
n^2 = gk\!\left(\frac{\rho_2-\rho_1}{\rho_2+\rho_1}-\frac{k^2T}{g(\rho_2+\rho_1)}\right)\!.
\end{equation}
According to equation~\eqref{eq:10-51}, \emph{if $\rho_2 < \rho_1$, the arrangement is stable; while if $\rho_2 > \rho_1$ the arrangement is unstable for all wave numbers in the range $0 < k < k_c$, where}
\begin{equation}
\label{eq:10-52}
k_c = \{(\rho_2-\rho_1)g/T\}^{1/2}.
\end{equation}
However, in the latter case, the arrangement is stable for all disturbances with $k > k_c$. Thus, surface tension succeeds in stabilizing a potentially unstable arrangement for all sufficiently short wavelengths; but the arrangement remains unstable for all sufficiently long wavelengths. Moreover, unlike when surface tension is absent, there is a \emph{mode of maximum instability} for which the amplitude of the disturbance grows most rapidly. The wave number $k_*$ of this most unstable mode is given by
\begin{equation}
\label{eq:10-53}
g(\rho_2-\rho_1) = 3k_*^2\,T,
\end{equation}
or
\begin{equation}
\label{eq:10-54}
k_* = k_c/\sqrt{3}.
\end{equation}
The corresponding value of $n$ is
\begin{equation}
\label{eq:10-55}
n_* = \left[\frac{2}{3^{3/2}}\frac{(\rho_2-\rho_1)g^2}{(\rho_2+\rho_1)\,T^{1/2}}\right]^{1/2}\!\!.
\end{equation}

\subsection*{(b) \textit{The case of exponentially varying density}}

A case of variable density for which a simple analytical solution can be found is
\begin{equation}
\label{eq:10-56}
\rho = \rho_0\,e^{\beta z},
\end{equation}
where $\beta$ is a constant. In this case equation~\eqref{eq:10-42} reduces to
\begin{equation}
\label{eq:10-57}
D^2w+\beta Dw - k^2(1-g\beta/n^2)w = 0.
\end{equation}
The general solution of this equation is
\begin{equation}
\label{eq:10-58}
w = A_1\,e^{q_1 z}+A_2\,e^{q_2 z},
\end{equation}
where $A_1$ and $A_2$ are two arbitrary constants and $q_1$ and $q_2$ are the roots of the equation
\begin{equation}
\label{eq:10-59}
q^2+q\beta - k^2(1-g\beta/n^2) = 0.
\end{equation}

If we suppose that the fluid is confined between two rigid planes at $z = 0$ and $z = d$, then the vanishing of $w$ at $z = 0$ is satisfied by the choice
\begin{equation}
\label{eq:10-60}
w = A_1(e^{q_1 z}-e^{q_2 z}),
\end{equation}
while the vanishing of $w$ at $z = d$ requires
\begin{equation}
\label{eq:10-61}
\exp[(q_1-q_2)d] = 1,
\end{equation}
or
\begin{equation}
\label{eq:10-62}
(q_1-q_2)d = 2im\pi,
\end{equation}
where $m$ is an integer.

Now writing the solution~\eqref{eq:10-60} in the manner
\begin{equation}
\label{eq:10-63}
w = A\,e^{\frac{1}{2}(q_1+q_2)z}[e^{\frac{1}{2}(q_1-q_2)z}-e^{-\frac{1}{2}(q_1-q_2)z}]
\end{equation}
and making use of~\eqref{eq:10-62} and the relation
\begin{equation}
\label{eq:10-64}
q_1+q_2 = -\beta
\end{equation}
(which follows from~\eqref{eq:10-59}), we have the alternative form
\begin{equation}
\label{eq:10-65}
w = \text{constant}\; e^{-\frac{1}{2}\beta z}\sin(m\pi z/d).
\end{equation}
Returning to equation~\eqref{eq:10-62} and inserting for $q_1$ and $q_2$ their values
\begin{equation}
\label{eq:10-66}
q_1 = \tfrac{1}{2}\{-\beta+\sqrt[\phantom{x}]{\beta^2+4k^2(1-g\beta/n^2)}\}
\end{equation}
and
\begin{equation}
\label{eq:10-67}
q_2 = \tfrac{1}{2}\{-\beta-\sqrt[\phantom{x}]{\beta^2+4k^2(1-g\beta/n^2)}\},
\end{equation}
we obtain
\begin{equation}
\label{eq:10-68}
\beta^2+4k^2\!\left(1-\frac{g\beta}{n^2}\right) = -\frac{4m^2\pi^2}{d^2},
\end{equation}
or
\begin{equation}
\label{eq:10-69}
\frac{g\beta}{n^2} = 1+\frac{\frac{1}{4}\beta^2 d^2+m^2\pi^2}{k^2 d^2}.
\end{equation}

From equation~\eqref{eq:10-68} it follows that \emph{the stratification is stable if $\beta$ is negative, while it is unstable if $\beta$ is positive}.

The smallest admissible value of $m$ is 1, which leads to the greatest numerical value of $n$. If $d$, $k$, and $m$ are given, $n^2$ is numerically greatest when $\beta$ is such that
\begin{equation}
\label{eq:10-70}
\tfrac{1}{4}\beta^2 d^2 = k^2d^2+m^2\pi^2.
\end{equation}

Rayleigh, who first worked out this example, has remarked: `Contrary to what is met with in most vibrating systems, there is (in the case of stability) a limit on the side of rapidity of vibration; but none on the side of slowness.'

\section{A general variational principle}\label{sec:93}

In \S\,92 we saw that in the absence of viscosity and surface tension the underlying characteristic value problem allows a variational formulation. We shall now show that such a formulation is possible in the general case.

The basic equations are (cf.\ equations~\eqref{eq:10-20} and~\eqref{eq:10-32})
\begin{align}
\label{eq:10-70a}
D\,\delta p &= -n\rho w + \mu(D^2-k^2)w + 2(D\mu)(Dw) + \frac{g}{n}(D\rho)w - \nonumber\\
&\quad -\frac{k^2}{n}\sum_s T_s\,w_s\,\delta(z-z_s) - \mu k^2 w + 2(D\mu)(Dw)
\end{align}
and
\begin{equation}
\label{eq:10-71}
k^2\delta p = -n\rho Dw - k^2\mu Dw + k^2wD\mu + D(\mu D^2w);
\end{equation}
and the boundary conditions with respect to which these equations must be solved are
\begin{equation}
\label{eq:10-72}
w \quad\text{and}\quad Dw\quad\text{vanish on a bounding surface.}
\end{equation}
The requirement that a solution of equations~\eqref{eq:10-70a} and~\eqref{eq:10-71} satisfy the necessary boundary conditions will lead to a determinate sequence of possible values for $n$. Let $n_i$ and $n_j$ denote two of these characteristic values; and let the solutions belonging to $n_i$ and $n_j$ be distinguished by the subscripts $i$ and $j$.

Now consider equation~\eqref{eq:10-70a} and~\eqref{eq:10-71} for the characteristic value $n_i$ and after multiplication by $w_j$ (belonging to $n_j$) integrate over the range of $z$ (which we shall assume to be $0 \leqslant z \leqslant d$). We have
\begin{align}
\label{eq:10-73}
\int w_j\,D\,\delta p_i\,dz &= \int\left\{-n_i\rho+\frac{g}{n_i}\,D\rho - k^2\mu\right\}w_iw_j\,dz - \nonumber\\
&\quad -\frac{k^2}{n_i}\sum_s T_s\,w_s(z_s)w_j(z_s) + \int w_j\{\mu D^2w_i+2(D\mu)(Dw_i)\}\,dz.
\end{align}
On integrating by parts the left-hand side of this equation, we have
\begin{equation}
\label{eq:10-74}
\int w_j\,D\,\delta p_i\,dz = -\int\delta p_i\,Dw_j\,dz,
\end{equation}
the integrated part vanishing on account of the boundary conditions. On the right-hand side of equation~\eqref{eq:10-74} we now substitute for $\delta p_i$ in accordance with~\eqref{eq:10-71}. We thus obtain
\begin{align}
\label{eq:10-75}
-\int\delta p_i\,Dw_j\,dz = \int\left\{\frac{n_i}{k^2}\rho+\mu\right\}(Dw_i)(Dw_j)\,dz - \frac{1}{k^2}\int(D\mu)w_i\,Dw_j\,dz \nonumber\\
{} - \int(D\mu)w_i\,Dw_j\,dz - \frac{1}{k^2}\int\mu(D^2w_i)(D^2w_j)\,dz.
\end{align}
The last integral on the right-hand side of this equation is
\begin{equation}
\label{eq:10-76}
\int(Dw_j)\,D(\mu D^2w_i)\,dz = -\int\mu(D^2w_i)(D^2w_j)\,dz,
\end{equation}
the integrated part again vanishing. We can now rewrite equation~\eqref{eq:10-75} in the form
\begin{align}
\label{eq:10-77}
-\int\delta p_i\,Dw_j\,dz &= \int\left\{\frac{n_i}{k^2}\rho+\mu\right\}(Dw_i)(Dw_j)\,dz + \nonumber\\
&\quad + \frac{1}{k^2}\int\mu(D^2w_i)(D^2w_j)\,dz - \int(D\mu)w_i(Dw_j)\,dz.
\end{align}

Now combining equations~\eqref{eq:10-73}, \eqref{eq:10-74}, and~\eqref{eq:10-77} we find after some further rearrangement
\begin{align}
\label{eq:10-78}
-n_i\int\rho\!\left\{w_iw_j+\frac{1}{k^2}(Dw_i)(Dw_j)\right\}\!dz + {} \nonumber\\
{} + \frac{g}{n_i}\int(D\rho)w_iw_j\,dz - \frac{k^2}{n_i}\sum_s T_s\,w_i(z_s)w_j(z_s) \nonumber\\
= \int\left\{\mu\!\left[k^2w_iw_j+(Dw_i)(Dw_j)+\frac{1}{k^2}(D^2w_i)(D^2w_j)\right]\right\}\!dz - \nonumber\\
{} - \int\left\{w_i[\mu D^2w_j + 2(D\mu)(Dw_j)] + w_j(D\mu)(Dw_j)\right\}\!dz.
\end{align}
Considering the last integral on the right-hand side of this equation, we can rearrange it in the manner
\begin{equation}
\label{eq:10-79}
-\int\{w_j\,D(\mu Dw_i)+(D\mu)D(w_iw_j)\}\,dz,
\end{equation}
or, after integrating by parts each of the two terms which occur under the integral sign, we are left with
\begin{equation}
\label{eq:10-80}
\int\{\mu(Dw_i)(Dw_j)+(D^2\mu)w_iw_j\}\,dz.
\end{equation}
Thus, we finally have
\begin{align}
\label{eq:10-81}
-n_i\int\rho\!\left\{w_iw_j+\frac{1}{k^2}(Dw_i)(Dw_j)\right\}\!dz + {} \nonumber\\
{} + \frac{g}{n_i}\int(D\rho)w_iw_j\,dz - \frac{k^2}{n_i}\sum_s T_s\,w_i(z_s)w_j(z_s) \nonumber\\
= \int\left\{(D^2\mu)w_iw_j + 2\mu(Dw_i)(Dw_j) + \frac{1}{k^2}\mu(D^2w_i)(D^2w_j)\right\}\!dz.
\end{align}

Interchanging $i$ and $j$ in equation~\eqref{eq:10-81} and subtracting the resulting equation from~\eqref{eq:10-81}, we obtain
\begin{align}
\label{eq:10-82}
(n_j-n_i)\!\left[\int\rho\!\left\{w_iw_j+\frac{1}{k^2}(Dw_i)(Dw_j)\right\}\!dz + {} \right. \nonumber\\
\left. {} + \frac{1}{n_in_j}\!\left\{g\int(D\rho)w_iw_j\,dz - k^2\sum_s T_s\,w_i(z_s)w_j(z_s)\right\}\right] = 0.
\end{align}

Hence, if $n_i \neq n_j$, we must have
\begin{align}
\label{eq:10-83}
\int\rho\!\left\{w_iw_j+\frac{1}{k^2}(Dw_i)(Dw_j)\right\}\!dz + {} \nonumber\\
{} + \frac{1}{n_in_j}\!\left\{g\int(D\rho)w_iw_j\,dz - k^2\sum_s T_s\,w_i(z_s)w_j(z_s)\right\} = 0 \quad (i\neq j).
\end{align}
A further relation of this kind can be obtained by rewriting equation~\eqref{eq:10-81} in the manner,
\begin{align}
\label{eq:10-84}
g\int(D\rho)w_iw_j\,dz - k^2\sum_s T_s\,w_i(z_s)w_j(z_s) \nonumber\\
= n_i^2\int\rho\!\left\{w_iw_j+\frac{1}{k^2}(Dw_i)(Dw_j)\right\}\!dz + n_i\int(D^2\mu)w_iw_j\,dz + \nonumber\\
{} + n_i\int\mu\!\left\{k^2w_iw_j + 2(Dw_i)(Dw_j) + \frac{1}{k^2}(D^2w_i)(D^2w_j)\right\}\!dz,
\end{align}
and treating it in the same way. We then obtain
\begin{align}
\label{eq:10-85}
(n_i+n_j)\int\rho\!\left\{w_iw_j+\frac{1}{k^2}(Dw_i)(Dw_j)\right\}\!dz + \int(D^2\mu)w_iw_j\,dz \nonumber\\
{} + \int\mu\!\left\{k^2w_iw_j + 2(Dw_i)(Dw_j) + \frac{1}{k^2}(D^2w_i)(D^2w_j)\right\}\!dz = 0 \quad (i\neq j).
\end{align}

Equations~\eqref{eq:10-83} and~\eqref{eq:10-85} enable us to draw certain general conclusions. Thus, if $n_i$ should be complex, we can suppose that $n_i$ and $n_j$ are complex conjugates and deduce from equation~\eqref{eq:10-83} that
\begin{equation}
\label{eq:10-86}
\int\rho\!\left\{|w|^2+\frac{1}{k^2}|Dw|^2\right\}\!dz + \frac{1}{|n|^2}\!\left\{g\int(D\rho)|w|^2\,dz - k^2\sum_s T_s|w(z_s)|^2\right\} = 0.
\end{equation}
In the absence of surface tension, equation~\eqref{eq:10-86} cannot be true if $Dp$ is everywhere positive so that in this case $n$ must be real.

Again, if $n$ should be complex we can deduce from equation~\eqref{eq:10-85} that
\begin{align}
\label{eq:10-87}
2\operatorname{re}(n)\int\rho\!\left\{|w|^2+\frac{1}{k^2}|Dw|^2\right\}\!dz \nonumber\\
= -\int(D^2\mu)|w|^2\,dz - \int\mu\!\left(k^2|w|^2+2|Dw|^2+\frac{1}{k^2}|D^2w|^2\right)\!dz.
\end{align}
From this equation it follows that \emph{if $n$ is complex, $\operatorname{re}(n) < 0$ if $D^2\mu$ is everywhere positive}. The restriction on $D^2\mu$ for the validity of this result is curious. Apart from this restriction, what the result implies is that if oscillatory modes exist they should be stable; conversely, overstability cannot occur.

\subsection*{(a) \textit{The variational principle}}

Returning to equation~\eqref{eq:10-81} and setting $i = j$, we get (on further suppressing the subscripts)
\begin{align}
\label{eq:10-88}
n\int\rho\!\left\{w^2+\frac{1}{k^2}(Dw)^2\right\}\!dz - \frac{g}{n}\int(D\rho)w^2\,dz + \frac{k^2}{n}\sum_s T_sw_s^2 \nonumber\\
= -\int\!\left\{\mu\!\left[k^2w^2+2(Dw)^2+\frac{1}{k^2}(D^2w)^2\right]+(D^2\mu)w^2\right\}\!dz.
\end{align}
This equation provides the basis for a variational formulation of the problem. To see this, consider the effect on $n$ (determined in accordance with~\eqref{eq:10-88}) of an arbitrary variation $\delta w$ in $w$ compatible only with the boundary conditions on $w$. We have to the first order
\begin{equation}
\label{eq:10-89}
-\left(I_1+\frac{1}{n^2}\left[gI_2-k^2\sum_s T_sw_s^2\right]\right)\!\delta n = n\delta I_1 - \frac{1}{n}[g\delta I_2 - 2k^2\sum_s T_sw_s\,\delta w_s] + \delta I_3,
\end{equation}
where
\begin{equation}
\label{eq:10-90}
I_1 = \int\rho\!\left\{w^2+\frac{1}{k^2}(Dw)^2\right\}\!dz, \quad I_2 = \int(D\rho)w^2\,dz,
\end{equation}
\begin{equation}
I_3 = \int\!\left\{\mu\!\left[k^2w^2+2(Dw)^2+\frac{1}{k^2}(D^2w)^2\right]+(D^2\mu)w^2\right\}\!dz,
\end{equation}
and $\delta I_1$, $\delta I_2$, and $\delta I_3$ are the corresponding variations in these integrals. After one or more integrations by parts, we find that these latter variations are given by
\begin{equation}
\label{eq:10-91}
\tfrac{1}{2}\delta I_1 = \int\delta w\!\left\{\rho w - \frac{1}{k^2}D(\rho Dw)\right\}\!dz,
\end{equation}
\begin{equation}
\label{eq:10-92}
\tfrac{1}{2}\delta I_2 = \int\delta w\{(D\rho)w\}\,dz,
\end{equation}
and
\begin{equation}
\label{eq:10-93}
\tfrac{1}{2}\delta I_3 = \int\delta w\!\left\{k^2\mu w - 2D(\mu Dw) + (D^2\mu)w + \frac{1}{k^2}D^2(\mu D^2w)\right\}\!dz.
\end{equation}

Combining the foregoing in accordance with~\eqref{eq:10-89}, we find after some further reductions that
\begin{align}
\label{eq:10-94}
\tfrac{1}{2}k^2\!\left\{I_1+\frac{1}{n^2}\left[gI_2-k^2\sum_s T_sw_s^2\right]\right\}\!\delta n \nonumber\\
= \int\delta w\!\left\{D\!\left[\rho Dw - \frac{\mu}{n}(D^2-k^2)Dw - \frac{1}{n}(D\mu)(D^2+k^2)w\right] - {}\right. \nonumber\\
{} - k^2\!\left[\rho w - \frac{\mu}{n}(D^2-k^2)w - \frac{g}{n^2}(D\rho)w - {}\right. \nonumber\\
\left.\left. {} + \frac{k^2}{n^2}\sum_s T_s\,\delta(z-z_s)w\right]\right\}\!dz.
\end{align}
We observe that the quantity which appears as a factor of $\delta w$ under the integral sign on the right-hand side of equation~\eqref{eq:10-94} vanishes if equation~\eqref{eq:10-41} governing $w$ is satisfied. Hence, a \emph{necessary and sufficient} condition for $\delta w$ to be zero to the first order for all small arbitrary variations in $w$ compatible with the boundary conditions is that $w$ be a solution of the characteristic value problem. A variational procedure of solving for the characteristic values is therefore possible.

\section{The case of two uniform viscous fluids separated by a horizontal boundary}\label{sec:94}

Let two uniform fluids of densities $\rho_1$ and $\rho_2$ and viscosities $\mu_1$ and $\mu_2$ be separated by a horizontal boundary at $z = 0$. The subscripts 1 and 2 distinguish the lower and the upper fluids, respectively.

In each of the two regions of constant $\rho$ and $\mu$, equation~\eqref{eq:10-41} reduces to
\begin{equation}
\label{eq:10-95}
D\!\left[\rho - \frac{\mu}{n}(D^2-k^2)\right]\!Dw = k^2\!\left[\rho - \frac{\mu}{n}(D^2-k^2)\right]\!w,
\end{equation}
where, for the present, we are suppressing the subscripts distinguishing the two fluids. Since $\rho$ and $\mu$ are constants, we can rewrite equation~\eqref{eq:10-95} in the form
\begin{equation}
\label{eq:10-96}
\left[1-\frac{\nu}{n}(D^2-k^2)\right](D^2-k^2)w = 0,
\end{equation}
where $\nu$ $(= \mu/\rho)$ is the coefficient of kinematic viscosity. The general solution of equation~\eqref{eq:10-96} is a linear combination of the solutions
\begin{equation}
\label{eq:10-97}
e^{\pm kz} \quad\text{and}\quad e^{\pm qz},
\end{equation}
where
\begin{equation}
\label{eq:10-98}
q^2 = k^2+n/\nu.
\end{equation}
Since $w$ must vanish both when $z \to -\infty$ (in the lower fluid) and $z \to +\infty$ (in the upper fluid), we can write
\begin{equation}
\label{eq:10-99}
w_1 = A_1\,e^{+kz}+B_1\,e^{+q_1 z} \quad (z < 0)
\end{equation}
and
\begin{equation}
\label{eq:10-100}
w_2 = A_2\,e^{-kz}+B_2\,e^{-q_2 z} \quad (z > 0),
\end{equation}
as the solutions appropriate for the two regions. In equations~\eqref{eq:10-99} and~\eqref{eq:10-100}, $A_1$, $B_1$, $A_2$, and $B_2$ are constants of integration and
\begin{equation}
\label{eq:10-101}
q_1 = \sqrt{(k^2+n/\nu_1)} \quad\text{and}\quad q_2 = \sqrt{(k^2+n/\nu_2)}.
\end{equation}
It should be noted that in writing the solutions for $w$ in the two regions $z < 0$ and $z > 0$ in the manner~\eqref{eq:10-99} and~\eqref{eq:10-100}, we have assumed that $q_1$ and $q_2$ are so defined that their real parts are positive.

On the interface $z = 0$, certain conditions must be satisfied; these have been found in \S\,91~(a) and given by~\eqref{eq:10-36} and~\eqref{eq:10-40}. The latter condition, under the circumstances of the present problem, gives
\begin{align}
\label{eq:10-102}
\left\{\!\left[\rho_1-\frac{\mu_1}{n}(D^2-k^2)\right]\!Dw_1\right\}_0 - \left\{\!\left[\rho_2-\frac{\mu_2}{n}(D^2-k^2)\right]\!Dw_2\right\}_0 \nonumber\\
= -\frac{k^2}{n^2}\,[g(\rho_2-\rho_1)-k^2T]w_{0s} - \frac{2k^2}{n}\,(\mu_1-\mu_2)(Dw)_0,
\end{align}
where $w_0$ and $(Dw)_0$ are the common values of $w_1$ and $w_2$, and, similarly, $Dw_1$ and $Dw_2$ at $z = 0$. Substituting for $w_1$ and $w_2$ in accordance with equations~\eqref{eq:10-99} and~\eqref{eq:10-100} in equation~\eqref{eq:10-102}, we find
\begin{equation}
\label{eq:10-103}
-k\rho_2A_2-k\rho_1A_1 = \frac{gk^2}{n^2}\!\left[(\rho_2-\rho_1)-\frac{k^2T}{g}\right]\!w_0 + \frac{k^2}{n}\,(\mu_1-\mu_2)(Dw)_0.
\end{equation}

Now applying the boundary conditions~\eqref{eq:10-36} and~\eqref{eq:10-103} to the solution given in~\eqref{eq:10-99} and~\eqref{eq:10-100} we obtain
\begin{gather}
A_1+B_1 = A_2+B_2 \quad (= w_0), \label{eq:10-104}\\
kA_1+q_1\,B_1 = -kA_2-q_2\,B_2 \quad (= Dw_0), \label{eq:10-105}\\
\mu_1(2k^2A_1+(q_1^2+k^2)B_1) = \mu_2(2k^2A_2+(q_2^2+k^2)B_2), \label{eq:10-106}
\end{gather}
and
\begin{align}
\label{eq:10-107}
-k\rho_2A_2-k\rho_1A_1 &= \frac{gk^2}{2n^2}(\rho_1-\rho_2) + \frac{k^2}{n}(A_1+B_1+A_2+B_2) + \nonumber\\
&\quad + \frac{k^2}{n}\,(\mu_1-\mu_2)(kA_1+q_1B_1-kA_2-q_2B_2).
\end{align}
Letting
\begin{equation}
\label{eq:10-108}
\alpha_1 = \frac{\rho_1}{\rho_1+\rho_2}, \quad \alpha_2 = \frac{\rho_2}{\rho_1+\rho_2} \quad (\alpha_1+\alpha_2 = 1),
\end{equation}
\begin{equation}
\label{eq:10-109}
R = \frac{gk|\rho_1-\rho_2|}{n^2|\rho_1+\rho_2|} + \frac{k^2T}{g(\rho_1+\rho_2)} = \frac{gk}{n^2}\!\left[(\alpha_1-\alpha_2)+\frac{k^2T}{g(\rho_1+\rho_2)}\right]\!,
\end{equation}
and
\begin{equation}
\label{eq:10-110}
C = \frac{k^2}{n}\frac{\mu_1-\mu_2}{\rho_1+\rho_2} = \frac{k^2}{n}(\alpha_1\nu_1-\alpha_2\nu_2),
\end{equation}
we can rewrite equations~\eqref{eq:10-104}--\eqref{eq:10-107}, in matrix notation, in the form of the single equation
\begin{equation}
\label{eq:10-111}
\begin{vmatrix}
1 & 1 & -1 & -1 \\[3pt]
k & q_1 & k & q_2 \\[3pt]
\dfrac{1}{2k^2\mu_1} & \mu_1(q_1^2+k^2) & -2k^2\mu_2 & -\mu_2(q_2^2+k^2) \\[6pt]
\frac{1}{2}R-C-\alpha_1 & \frac{1}{2}R-q_1C/k & \frac{1}{2}R+C-\alpha_2 & \frac{1}{2}R-\alpha_2q_2C/k\,q_1
\end{vmatrix}\!\begin{pmatrix} A_1 \\ B_1 \\ A_2 \\ B_2 \end{pmatrix} = 0.
\end{equation}
The determinant of the linear system of equations which~\eqref{eq:10-111} represents must clearly vanish. The determinant can be reduced by subtracting the first column from the second and similarly the third column from the fourth. By this procedure and by making further use of equations~\eqref{eq:10-98} and~\eqref{eq:10-108}, we obtain
\begin{equation}
\label{eq:10-112}
\begin{vmatrix}
q_1-k & 2k & q_2-k \\
\alpha_1(q_1^2-\alpha_1q_2\nu_2) & 2k(\alpha_1\nu_1-\alpha_2\nu_2) & -\alpha_2 n \\
\alpha_1+C(1-q_1/k) & R-1 & \alpha_2-C(1-q_2/k)
\end{vmatrix} = 0.
\end{equation}
On expanding this determinant and simplifying, we find
\begin{align}
\label{eq:10-113}
-\left[\frac{gk}{n^2}\!\left[(\alpha_1-\alpha_2)+\frac{k^2T}{g(\rho_1+\rho_2)}\right]+1\right](\alpha_1q_1+\alpha_2q_2-k) - 4k\alpha_1\alpha_2 + {} \nonumber\\
{} + \frac{4k^2}{n^2}(\alpha_1\nu_1-\alpha_2\nu_2)(\alpha_1q_1-\alpha_2q_2)+k(\alpha_1-\alpha_2)^2 + {} \nonumber\\
{} + \frac{4k^2}{n^2}(\alpha_1\nu_1-\alpha_2\nu_2)^2(q_1-k)(q_2-k) = 0,
\end{align}
where it may be recalled that $q_1$ and $q_2$ are related to $n$ by equations~\eqref{eq:10-101}. Equation~\eqref{eq:10-113} is the required characteristic equation for $n$: it was first derived, essentially in this form, by Harrison in 1908.

\subsection*{(a) \textit{The case $\nu_1 = \nu_2$}}

In the further discussion of equation~\eqref{eq:10-113} we shall restrict ourselves to the case when the kinematic viscosities of the two fluids are the same, i.e.\ when $\nu_1 = \nu_2$. This assumption simplifies the characteristic equation~\eqref{eq:10-113} considerably; but one would not expect that any of the essential features of the problem would be obscured by this simplifying assumption. When $\nu_1 = \nu_2$,
\begin{equation}
\label{eq:10-114}
q_1 = q_2 = k\sqrt{(1+n/k^2\nu)} = q \quad\text{(say),}
\end{equation}
and equation~\eqref{eq:10-113} reduces to
\begin{align}
\label{eq:10-115}
-\left[\frac{gk}{n^2}\!\left[(\alpha_1-\alpha_2)+\frac{k^2T}{g(\rho_1+\rho_2)}\right]+1\right](q-k) - \frac{4k\nu}{n}(\alpha_1-\alpha_2)^2(q-k) + {} \nonumber\\
{} + \frac{4k^2\nu^2}{n^2}(\alpha_1-\alpha_2)^2(q-k)^2 = 0.
\end{align}
Letting
\begin{equation}
z = n/k^2\nu. \label{eq:10-117}
\end{equation}
so that
\begin{equation}
q = k\sqrt{(1+z)},
\end{equation}
we can rewrite equation~\eqref{eq:10-115} in the form
\begin{align}
\label{eq:10-118}
\frac{g}{k^3\nu^2}\!\left[(\alpha_1-\alpha_2)+\frac{k^2T}{g(\rho_1+\rho_2)}\right]\frac{1}{z^2}\cdot\frac{1+\frac{4}{z}(\alpha_1-\alpha_2)^2}{[\sqrt{(1+z)}-1]} + \frac{4\alpha_1\alpha_2}{z^2}\cdot {} \nonumber\\
{} - \frac{4}{z}(\alpha_1-\alpha_2)\sqrt{(1+z)}\cdot[(1+z)-1] + \frac{4\alpha_1\alpha_2}{z^2} = 0.
\end{align}
With the further substitutions
\begin{equation}
\label{eq:10-119}
y = q/k = \sqrt{(1+z)}, \quad x = y^2-1,
\end{equation}
\begin{equation}
\label{eq:10-120}
Q = \frac{g}{k^3\nu^2}, \quad S = \frac{T}{(\rho_1+\rho_2)(\nu\nu)^{4/3}}, \quad\text{and}\quad Q^{4/3}S = \frac{k^2T}{g(\rho_1+\rho_2)},
\end{equation}
we find that equation~\eqref{eq:10-118} reduces to the following quartic for $y$:
\begin{align}
\label{eq:10-121}
y^4+4\alpha_1\alpha_2\,y^3+2(1-6\alpha_1\alpha_2)y^2-4(1-3\alpha_1\alpha_2)y + {} \nonumber\\
{} + (1-4\alpha_1\alpha_2)+Q(\alpha_1-\alpha_2)+Q^{4/3}S = 0.
\end{align}

Since in expressing the solutions for $w_1$ and $w_2$ in the manner~\eqref{eq:10-99} and~\eqref{eq:10-100} we have supposed that the real parts of $q_1$ and $q_2$ are positive, it is clear that in the present case (cf.\ equations~\eqref{eq:10-117} and~\eqref{eq:10-119}) only those roots of equation~\eqref{eq:10-121} which have positive real parts lead to physically valid solutions.

For assigned $\alpha_1$ $(= 1-\alpha_2)$, $Q$, and $S$, a root of the quartic~\eqref{eq:10-121} whose real part is positive determines a pair of values of $k$ and $n$ which belong to each other; for, according to equations~\eqref{eq:10-116},~\eqref{eq:10-119}, and~\eqref{eq:10-120},
\begin{equation}
\label{eq:10-122}
k = \left(\frac{g}{Q}\right)^{1/3}\!\!\frac{1}{Q^{1/3}} \quad\text{and}\quad n = k^2\nu(y^2-1) = \left(\frac{g}{Q}\right)^{2/3}\!\!\frac{y^2-1}{Q^{4/3}}.
\end{equation}
Hence if $k$ and $n$ are measured in the units
\begin{equation}
\label{eq:10-123}
(g/\nu^2)^{1/3}\quad\text{and}\quad (g^2/\nu)^{1/3}\;\text{sec}^{-1},\;\text{respectively,}
\end{equation}
we can write
\begin{equation}
\label{eq:10-124}
k = Q^{-1/3} \quad\text{and}\quad n = (y^2-1)Q^{-2/3}.
\end{equation}

\subsection*{(b) \textit{The modes of maximum instability for the case $\nu_1 = \nu_2$, $\rho_2 > \rho_1$, and $S = 0$}}

When $S = 0$, the equation we have to consider is
\begin{align}
\label{eq:10-125}
y^4+4\alpha_1\alpha_2\,y^3+2(1-6\alpha_1\alpha_2)y^2-4(1-3\alpha_1\alpha_2)y + {} \nonumber\\
{} + (1-4\alpha_1\alpha_2)+Q(\alpha_1-\alpha_2) = 0.
\end{align}
It can be readily shown that this equation, in the case when $\alpha_2 > \alpha_1$, allows only one root whose real part is positive. Indeed, this root is real and, moreover, $y > 1$. Accordingly, $n$ is real and positive (cf.\ equation~\eqref{eq:10-124}) and the amplitude of the disturbance will grow exponentially with time. \emph{The arrangement is therefore unstable for disturbances of all wave numbers.}

The asymptotic behaviour of $n$ for $k \to 0$ and $k \to \infty$ can be derived by considering equation~\eqref{eq:10-125} for $y \to \infty$ and $y \to 1$. In the former case
\begin{equation}
\label{eq:10-126}
Q \to y^4/(\alpha_2-\alpha_1) \quad (y \to \infty),
\end{equation}
and we deduce from~\eqref{eq:10-124} that simultaneously
\begin{equation}
\label{eq:10-127}
k \to y^{-4/3}(\alpha_2-\alpha_1)^{1/3} \quad\text{and}\quad n \to y^{-2/3}(\alpha_2-\alpha_1)^{2/3}
\end{equation}
so that
\begin{equation}
\label{eq:10-128}
n^3 \to (\alpha_2-\alpha_1)k \quad (k \to 0).
\end{equation}
This asymptotic relation for $k \to 0$ is exactly that which obtains in the absence of viscosity. This agreement is what one should expect on general grounds, namely, that viscosity plays no role among the very long wavelengths.

Considering next the limit $y \to 1$, we have (see equation~\eqref{eq:10-132} below)
\begin{equation}
\label{eq:10-129}
Q \to 4(y-1)/(\alpha_2-\alpha_1) \quad (y \to 1),
\end{equation}
and we deduce that simultaneously
\begin{equation}
\label{eq:10-130}
k \to \left(\frac{\alpha_2-\alpha_1}{4}\right)^{1/3}\!(y-1)^{-1/3} \quad\text{and}\quad n \to 2\!\left(\frac{\alpha_2-\alpha_1}{4}\right)^{2/3}\!(y-1)^{1/3}
\end{equation}
so that
\begin{equation}
\label{eq:10-131}
n \to (\alpha_2-\alpha_1)/2k \quad (k \to \infty).
\end{equation}
According to equations~\eqref{eq:10-128} and~\eqref{eq:10-131}, $n \to 0$ both when $k \to 0$ and $k \to \infty$. There exists, therefore, a \emph{mode of maximum instability}. And this can be accomplished most easily in the following manner.

Rewriting equation~\eqref{eq:10-125} in the form
\begin{align}
\label{eq:10-132}
Q &= \frac{1}{\alpha_2-\alpha_1}\{y^4+4\alpha_1\alpha_2\,y^3+2(1-6\alpha_1\alpha_2)y^2 - 4(1-3\alpha_1\alpha_2)y + {} \nonumber\\
&\quad {}+ (1-4\alpha_1\alpha_2)\} \nonumber\\
&= \frac{y-1}{\alpha_2-\alpha_1}\{y^3+(1+4\alpha_1\alpha_2)y^2+(3-8\alpha_1\alpha_2)y - (1-4\alpha_1\alpha_2)\}
\end{align}
and evaluating $Q$ for various assigned values of $y \geqslant 1$, we can first obtain corresponding pairs of values of $Q$ and $y$; then, making use of equations~\eqref{eq:10-124}, we can convert them into corresponding pairs of $k$ and $n$. The $(k,n)$-relationships derived in this manner are illustrated in Figs.~106 and 107; and the modes of maximum instability for various values of $\alpha_2-\alpha_1 = (\rho_2-\rho_1)/(\rho_1+\rho_2)$ are listed in Table~XLVI.

\figurePlaceholder{106}{The dependence of the rate of growth $n$ (measured in the unit $(g^2/\nu)^{1/3}$) of a disturbance on its wave number $k$ (measured in the unit $(g/\nu^2)^{1/3}$) in case the upper fluid is more dense and the kinematic viscosities are the same. The curves labelled 1, 2, 3, and 4 are for values of $(\rho_2-\rho_1)/(\rho_1+\rho_2)$ $(= \alpha_2-\alpha_1)$ = 0.01, 0.05, 0.10, and 0.15, respectively.}

\figurePlaceholder{107}{Same as Fig.~106 but for values $\alpha_2-\alpha_1 = 0.25$, 0.5, 0.9, and 1.0 (curves labelled 5, 6, 7, and 8, respectively).}

\begin{table}[ht]
\centering
\caption{The modes of maximum instability for the case $\nu_1 = \nu_2$ and $\rho_2 > \rho_1$}
\label{tab:XLVI}
\small
\begin{tabular}{ccccc}
\hline
$\dfrac{\rho_2-\rho_1}{\rho_1+\rho_2}$ & $k(g/\nu^2)^{-1/3}$ & $n(\nu/g^2)^{1/3}$ \\[6pt]
\hline
0.01 & 0.1134 & 0.02081 \\
0.05 & 0.1442 & 0.05686 \\
0.10 & 0.1793 & 0.09063 \\
0.15 & 0.3394 & 0.1252 \\
0.25 & 0.3806 & 0.2342 \\
0.50 & 0.4112 & 0.4265 \\
1.00 & 0.4907 & 0.4590 \\
\hline
\end{tabular}
\end{table}

\subsection*{(c) \textit{The effect of surface tension on the unstable modes for $\nu_1 = \nu_2$ and $\rho_2 > \rho_1$}}

When surface tension is taken into account, equation~\eqref{eq:10-132} is replaced by
\begin{equation}
\label{eq:10-133}
Q^{4/3}\!\left(Q^{1/3}-\frac{S}{\alpha_2-\alpha_1}\right) = \frac{y-1}{(\alpha_2-\alpha_1)}\{y^3+(1+4\alpha_1\alpha_2)y^2+(3-8\alpha_1\alpha_2)y-(1-4\alpha_1\alpha_2)\}.
\end{equation}
From this equation it follows that for a given $S$, as $y$ goes from $1$ to $\infty$, the left-hand side takes all values from 0 to $\infty$. The left-hand side is zero when
\begin{equation}
\label{eq:10-134}
Q^{1/3} = S/(\alpha_2-\alpha_1) = Q_c^{1/3} \quad\text{(say);}
\end{equation}
so that $y = 1$ corresponds to $Q = Q_c$. This means that unstable modes occur only for $k < k_c$ where
\begin{equation}
\label{eq:10-135}
k_c = Q_c^{-1/3} = \sqrt{[(\alpha_2-\alpha_1)/S]}.
\end{equation}
It can be readily verified that this wave number $k_c$ (now measured in the unit specified in~\eqref{eq:10-123}) is the same as~\eqref{eq:10-52} which we found in the absence of viscosity. \emph{The wave numbers which are stabilized by surface tension are, therefore, independent of viscosity.} (This result does not, in fact, depend on the assumption $\nu_1 = \nu_2$ which underlies the present derivation.)

When $y \to \infty$, $Q \to \infty$ and the term in $S$ on the left-hand side of~\eqref{eq:10-133} can be ignored. This means that for $k \to 0$ the same asymptotic behaviour~\eqref{eq:10-128} obtains as in the absence of surface tension.

For the particular case $\alpha_1 = 0$, the $(n, k)$-relationships, which have been derived by Reid for a number of values of $S$, are illustrated in Fig.~108. And in Table~XLVII we present his results on the dependence of the modes of maximum instability on $S$.

\begin{table}[ht]
\centering
\caption{The modes of maximum instability for the case $\nu_1 = \nu_2$ and $\rho_1 = 0$ and different values of $S$}
\label{tab:XLVII}
\small
\begin{tabular}{ccc}
\hline
$S$ & $k(g/\nu^2)^{-1/3}$ & $n(\nu/g^2)^{1/3}$ \\
\hline
0 & 0.4997 & 0.4599 \\
$\frac{1}{4}$ & 0.4591 & 0.4524 \\
$\frac{1}{2}$ & 0.4227 & 0.4443 \\
1 & 0.4157 & 0.4315 \\
\hline
\end{tabular}
\end{table}

In addition to the unstable modes we have considered, the problem also allows damped modes when $S \neq 0$. These damped modes exhibit a curious behaviour. For example, Reid has shown that in the case $\alpha_1 = 0$, the damped modes are all aperiodic so long as $S < 1.3165$; while for $S$ greater than this critical value, damped oscillatory modes become possible for certain ranges of the wave number.

\figurePlaceholder{108}{The dependence of rate of growth of a disturbance on its wave number and on the surface tension parameter $S$, in case the lower fluid is of zero density; $n$ and $k$ are measured in the same units as in Figs.~106 and 107.}

\subsection*{(d) \textit{The manner of decay in the case, $\nu_1 = \nu_2$, $\rho_2 < \rho_1$, and $S = 0$}}

When $\alpha_2 < \alpha_1$, equation~\eqref{eq:10-125} admits two roots whose real parts are positive. More precisely, the behaviour of these roots for assigned values of $Q$ is as follows.

When $\alpha_2 < \alpha_1$, equation~\eqref{eq:10-132} gives positive values of $Q$ for positive real values of $y$ only in a range, $y_a \leqslant y < 1$, where $y_a > 0$ is a determinate constant depending only on $\alpha_2$ and $\alpha_a$ $(= 1-\alpha_2)$. Further, in the range, $y_a \leqslant y \leqslant 1$, $Q$ attains a maximum (say $Q_a$) and vanishes at both ends. This means that there exists a value of $k$ (say $k_a$) such that for $k > k_a$ $(= Q_a^{-1/3})$ there are two possible modes of decay: a `viscous' mode which decays very rapidly and a `creeping' mode which decays very slowly. The difference between these two modes becomes increasingly pronounced as $k \to \infty$. Both of these modes decay aperiodically. (We shall return, presently, to the difference in the proper functions belonging to the two modes.) For $Q > Q_a$ equation~\eqref{eq:10-125} admits two complex roots which are conjugates of one another and whose real parts

\figurePlaceholder{109}{The dependence of the rate of decay $-\operatorname{re}(n)$ (measured in the unit $(g^2/\nu)^{1/3}$) of a disturbance on the wave number $k$ (measured in the unit $(g/\nu^2)^{1/3}$) in case the lower fluid is more dense and the kinematic viscosities are the same. The curves labelled 1 and 2 are for values of $\alpha_2-\alpha_1 = 0.01$ and 0.10, respectively. The dashed parts of each curve refer to the rapidly decaying `viscosity' mode.}

are positive. For $k < k_a$, the damping, therefore, takes place as oscillations of decreasing amplitude.

The establishment of the $(k,n)$-relationships for $k > k_a$ can be accomplished most simply by evaluating $Q$ (according to equation~\eqref{eq:10-132}) for various values of $y$ in the range, $y_a \leqslant y \leqslant 1$, and converting the resulting pairs of $Q$ and $y$ into corresponding pairs of $k$ and $n$ in accordance with~\eqref{eq:10-124}. For $Q > Q_a$ there would appear to be no alternative to solving the quartic equation~\eqref{eq:10-125} explicitly.

The $(k,n)$-relationships derived in the manner we have described are illustrated in Figs.~109, 110, and 111.

\figurePlaceholder{110}{Same as Fig.~109 but for values of $\alpha_2-\alpha_1 = 0.05$, 0.15, 0.25, 0.50, and 0.90 (curves labelled 2, 4, 5, 6, and 7, respectively).}

\figurePlaceholder{111}{The dependence of the imaginary part, $\operatorname{im}(n)$, of $n$ on the wave number $k$ in case the lower fluid is more dense. The curves labelled 1, 2, 3, 4, 5, 6, 7, and 8 are for values of $\alpha_2-\alpha_1 = 0.01$, 0.05, 0.10, 0.15, 0.25, 0.50, 0.90, and 1.0, respectively.}

this means that there is a maximum to the frequency of oscillations which can be set up or maintained.

In Table~XLVIII the branch points $(k_a, n)$ at which the manner of the decay changes its character from being aperiodic to being periodic are listed for various values of $\alpha_1-\alpha_2$.

As we have already remarked, the two aperiodic modes of decay which are possible for $k > k_a$ are quite different in that one of them decays very slowly relative to the other. The physical origin for this difference in the two modes can be understood in terms of the proper solutions belonging to them. In Fig.~112 the solutions for the two modes (normalized to unit maximum amplitude) which occur for $k = 1.1$ and $\alpha_1-\alpha_2 = 0.5$ are illustrated. The two modes are characterized by $n = -0.2757$ and $-1.0439$. We observe that the solution belonging to the more rapidly decaying viscous mode decreases in amplitude much more slowly than the solution belonging to the less rapidly decaying creeping mode: the origin of the greater damping of the viscous mode is thus made apparent. For comparison we have also included in Fig.~112 the proper solution for the oscillatory mode (characterized by $n = -0.1947 \pm 0.3523\,i$) which occurs for $k = 0.5$.

\begin{table}[ht]
\centering
\caption{The branch point $(k_a, n)$ at which the manner of decay changes from being aperiodic to periodic}
\label{tab:XLVIII}
\small
\begin{tabular}{ccc}
\hline
$\dfrac{\rho_1-\rho_2}{\rho_1+\rho_2}$ & $k_a(g/\nu^2)^{1/3}$ & $n(\nu/g^2)^{1/3}$ \\[6pt]
\hline
0.01 & 0.2537 & $-0.0797$ \\
0.05 & 0.4321 & $-0.1101$ \\
0.10 & 0.3446 & $-0.1747$ \\
0.15 & 0.6436 & $-0.1285$ \\
0.25 & 0.7460 & $-0.3200$ \\
0.50 & 0.9355 & $-0.5028$ \\
0.90 & 1.0324 & $-0.5215$ \\
1.00 & 1.1537 & $-0.7673$ \\
\hline
\end{tabular}
\end{table}

\figurePlaceholder{112}{(a) The proper solutions for the creeping $c$ and the viscous $v$ modes for the case $\alpha_1-\alpha_2 = 0.5$ and $k = 1.1$; the decay rates associated with these modes are, respectively, $n = -0.2757$ and $n = -1.0439$. (b) The proper solutions for the oscillatory mode for the case $\alpha_1-\alpha_2 = 0.5$ and $k = 0.5$; the complex value of $n$ associated with this mode is $n = -0.1947 \pm 0.3523\,i$.}

\subsection*{(e) \textit{Gravity waves}}

An interesting by-product of the preceding analysis is the dispersion relation for the so-called gravity waves which occur in an infinitely deep ocean. We can obtain this relation by passing to the limit $\alpha_2 = 0$ and $\alpha_1 = 1$. In this limit, the characteristic equation becomes (cf.\ equation~\eqref{eq:10-121})
\begin{equation}
\label{eq:10-136}
y^4+2y^3-4y+1+Q+Q^{4/3}S = 0.
\end{equation}
The $(k,n)$-relationships derived from this equation (in the manner described in \S\,(c) above) for several values of $S$ are shown in Fig.~113.

\figurePlaceholder{113}{The dependence of the real (a) and the imaginary (b) parts of $n$ on the wave number $k$ and the surface tension parameter $S$ for gravity waves.}

And in Table~XLIX the coordinates of the branch point (where the manner of the damping changes) are listed. The proper solutions belonging to $k = 0.8$ and 1.4 are shown in Fig.~114; the characteristic difference in the solutions belonging to the viscous and the creeping modes is clearly exhibited by the curves for $k = 1.4$.

\begin{table}[ht]
\centering
\caption{The coordinates of the branch point at which the manner of decay changes from being aperiodic to periodic for gravity waves}
\label{tab:XLIX}
\small
\begin{tabular}{ccc}
\hline
$S$ & $k_a(g/\nu^2)^{1/3}$ & $n(\nu/g^2)^{1/3}$ \\
\hline
0 & 1.1061 & $-0.7953$ \\
$\frac{1}{2}$ & 1.1053 & $-0.8557$ \\
$\frac{1}{4}$ & 1.3600 & $-0.9884$ \\
1 & 1.7935 & $-1.3085$ \\
2 & 1.8131 & $-1.7569$ \\
4 & 2.1073 & $-3.3732$ \\
\hline
\end{tabular}
\end{table}

\figurePlaceholder{114}{(a) The proper solutions for the creeping $c$ and the viscous $v$ modes for gravity waves of wave number $k = 1.4$. The damping constants associated with these modes are $n = -0.4314$ and $n = -1.6529$, respectively. (b) The proper solutions for gravity waves of wave number $k = 0.8$; for this value of $k$ the waves are periodically damped with $n = -0.4404 \pm 0.6061\,i$.}

\section{The effect of rotation}\label{sec:95}

We shall consider the effect of rotation on the equilibrium of a stratified heterogeneous inviscid fluid. The initial kinematic state is, then, the same as that postulated in \S\,91 except that now the fluid will be supposed to be inviscid but subject to a uniform rotation with an angular velocity $\Omega$ about the vertical. Under these circumstances the relevant perturbation equations are (cf.\ Chapter III, \S\S\,21 and 22):
\begin{gather}
\rho\frac{\partial u}{\partial t} - 2\rho\Omega v = -\frac{\partial}{\partial x}\,\delta p, \label{eq:10-137}\\
\rho\frac{\partial v}{\partial t} + 2\rho\Omega u = -\frac{\partial}{\partial y}\,\delta p, \label{eq:10-138}\\
\rho\frac{\partial w}{\partial t} = -\frac{\partial}{\partial z}\,\delta p - g\,\delta\rho + \sum_s T_s\!\left(\frac{\partial^2}{\partial x^2}+\frac{\partial^2}{\partial y^2}\right)\!\delta z_s\,\delta(z-z_s), \label{eq:10-139}\\
\delta\rho = -wDp, \label{eq:10-140}\\
\frac{\partial u}{\partial x}+\frac{\partial v}{\partial y}+\frac{\partial w}{\partial z} = 0, \label{eq:10-141}
\end{gather}
where in equation~\eqref{eq:10-139} we have allowed for the possibility that at some preassigned levels $z_s$, the density may change discontinuously and bring into play effects due to surface tension.

For solutions having the $(x,y,t)$-dependence given by~\eqref{eq:10-13}, equations~\eqref{eq:10-137}--\eqref{eq:10-141} become
\begin{gather}
pnu - 2\rho\Omega v = -ik_x\,\delta p, \label{eq:10-142}\\
pnv + 2\rho\Omega u = -ik_y\,\delta p, \label{eq:10-143}\\
pnw = -D\,\delta p - g\,\delta\rho - k^2\sum_s T_s\,\delta z_s\,\delta(z-z_s), \label{eq:10-144}\\
n\,\delta\rho = -wDp, \label{eq:10-145}\\
ik_x\,u + ik_y\,v = -Dw. \label{eq:10-146}
\end{gather}
Multiplying equations~\eqref{eq:10-142} and~\eqref{eq:10-143} by $-ik_x$ and $-ik_y$, respectively, adding, and making use of equation~\eqref{eq:10-146}, we obtain
\begin{equation}
\label{eq:10-147}
\rho nDw + 2\rho\Omega\zeta = -k^2\delta p,
\end{equation}
where
\begin{equation}
\label{eq:10-148}
\zeta = ik_yv - ik_xu
\end{equation}
is the $z$-component of the vorticity. Similarly, by multiplying equations~\eqref{eq:10-142} and~\eqref{eq:10-143} by $-ik_y$ and $+ik_x$, respectively, adding, and making use of equations~\eqref{eq:10-146} and~\eqref{eq:10-148}, we obtain
\begin{equation}
\label{eq:10-149}
\rho n\zeta - 2\rho\Omega Dw = 0
\end{equation}
or
\begin{equation}
\label{eq:10-150}
\zeta = 2\Omega Dw/n.
\end{equation}
Inserting this expression for $\zeta$ in equation~\eqref{eq:10-147}, we find
\begin{equation}
\label{eq:10-151}
\rho n\!\left(1+\frac{4\Omega^2}{n^2}\right)\!Dw = -k^2\delta p;
\end{equation}
whereas by combining equations~\eqref{eq:10-144} and~\eqref{eq:10-145} and making use, also, of the relation (cf.\ equation~\eqref{eq:10-31}),
\begin{equation}
\label{eq:10-152}
\delta z_s = w_s/n \quad (w_s = w(z_s)),
\end{equation}
we find
\begin{equation}
\label{eq:10-153}
D\,\delta p = -\rho nw + \frac{g}{n}(D\rho)w - \frac{k^2}{n}\sum_s T_s\,w_s\,\delta(z-z_s).
\end{equation}
Now eliminating $\delta p$ between equations~\eqref{eq:10-151} and~\eqref{eq:10-153} we obtain
\begin{equation}
\label{eq:10-154}
\left(1+\frac{4\Omega^2}{n^2}\right)D(\rho Dw) - k^2\rho w = -\frac{gk^2}{n^2}(D\rho)w - \frac{k^4}{n^2}\sum_s T_s\,\delta(z-z_s)w.
\end{equation}
Letting
\begin{equation}
\label{eq:10-155}
\kappa^2 = \frac{k^2}{1+4\Omega^2/n^2},
\end{equation}
we can rewrite equation~\eqref{eq:10-154} in the form
\begin{equation}
\label{eq:10-156}
D(\rho Dw) - \kappa^2\rho w = -\frac{g\kappa^2}{n^2}\left\{(D\rho) - \frac{k^2}{g}\sum_s T_s\,\delta(z-z_s)\right\}w.
\end{equation}
(We shall presently see that either $n^2 > 0$ or $n^2 < -4\Omega^2$ so that $\kappa^2$ is always positive and the definition~\eqref{eq:10-155} a permissible one.)

At the interface $z_s$ between two fluids, equation~\eqref{eq:10-156} leads to the condition (cf.\ equation~\eqref{eq:10-49})
\begin{equation}
\label{eq:10-157}
\Delta_s(\rho Dw) = -\frac{\kappa^2}{n^2}\,[g\,\Delta_s(\rho)-k^2T_s]w_s.
\end{equation}
While
\begin{equation}
\label{eq:10-158}
D(\rho Dw) - \kappa^2\rho w = -\frac{g\kappa^2}{n^2}(D\rho)w \quad (z \neq z_s).
\end{equation}
We observe that equation~\eqref{eq:10-158} is formally identical with equation~\eqref{eq:10-42} which is valid in the absence of rotation; the difference in the two problems arises only from the circumstance that we now have $\kappa^2$ in place of $k^2$.

\subsection*{(a) \textit{The case of two uniform fluids separated by a horizontal boundary}}

In view of the formal identity of equations~\eqref{eq:10-42} and~\eqref{eq:10-158}, we can write down directly the solution for the present problem from the solution for the same problem, in the absence of rotation, given in \S\,92\,(a). Thus, in place of the solution~\eqref{eq:10-51}, we now have
\begin{equation}
\label{eq:10-159}
n^2 = g\kappa\!\left[\frac{\rho_2-\rho_1}{\rho_2+\rho_1}-\frac{k^2T}{g(\rho_1+\rho_2)}\right]\!,
\end{equation}
or, substituting for $\kappa$ its value~\eqref{eq:10-155}, we have
\begin{equation}
\label{eq:10-160}
n^2\!\left(1+\frac{4\Omega^2}{n^2}\right)^{1/2} = gk\!\left[\frac{\rho_2-\rho_1}{\rho_2+\rho_1}-\frac{k^2T}{g(\rho_1+\rho_2)}\right]\!.
\end{equation}

In terms of $n_0$, the value of $n$ appropriate for the wave number $k$ in the absence of rotation, we can rewrite equation~\eqref{eq:10-160} in the manner,
\begin{equation}
\label{eq:10-161}
n^2\sqrt{(1+4\Omega^2/n^2)} = n_0^2.
\end{equation}
The solution of this equation is
\begin{equation}
\label{eq:10-162}
n^2 = -2\Omega^2+\sqrt{(4\Omega^2+n_0^4)} \quad\text{if}\; n_0^2 > 0
\end{equation}
and
\begin{equation}
\label{eq:10-163}
n^2 = -2\Omega^2-\sqrt{(4\Omega^2+n_0^4)} \quad\text{if}\; n_0^2 < 0.
\end{equation}
From equations~\eqref{eq:10-162} and~\eqref{eq:10-163} it follows that in the present case rotation does not affect the instability or stability, as such, of a stratification: $n^2 > 0$ if $n_0^2 > 0$ and $n^2 < 0$ if $n_0^2 < 0$. On the other hand, according to equation~\eqref{eq:10-163}
\begin{equation}
\label{eq:10-164}
n^2 \to -4\Omega^2 \quad\text{if}\; n_0^2 \to 0\;\text{through negative values.}
\end{equation}
\emph{The frequency of the stable oscillations cannot, thus, be less than the natural frequency of wave propagation in a rotating fluid} (cf.\ \S\,23, equation III~(71)). It appears that this is a general result.

\subsection*{(b) \textit{The case of exponentially varying density}}

From the solution for the same problem in the absence of rotation given in \S\,92\,(b) (equation~\eqref{eq:10-69}), we can now write down
\begin{equation}
\label{eq:10-165}
\frac{g\beta}{n^2} = 1+\frac{\frac{1}{4}\beta^2d^2+m^2\pi^2}{\kappa^2 d^2},
\end{equation}
or, substituting for $\kappa$ its value~\eqref{eq:10-155}, we have
\begin{equation}
\label{eq:10-166}
\frac{g\beta}{n^2} = 1+\frac{\frac{1}{4}\beta^2d^2+m^2\pi^2}{k^2d^2}\!\left(1+\frac{4\Omega^2}{n^2}\right)\!.
\end{equation}
From this equation we obtain
\begin{equation}
\label{eq:10-167}
n^2 = n_0^2\!\left(1-\frac{4\Omega^2}{g\beta}\cdot\frac{\frac{1}{4}\beta^2d^2+m^2\pi^2}{k^2d^2}\right)\!,
\end{equation}
where $n_0^2$ is the value of $n^2$ in the absence of rotation given by~\eqref{eq:10-68}. Remembering that $n_0^2$ is positive or negative, according as $\beta$ is positive or negative, we may conclude from~\eqref{eq:10-167} that if $\beta$ is negative the stratification is stable for disturbances of all wave numbers, and moreover,
\begin{equation}
\label{eq:10-168}
n^2 \to n_0^2\;\text{for}\;k\to\infty \quad\text{and}\quad n^2 \to -4\Omega^2\;\text{for}\;k\to 0;
\end{equation}
whereas if $\beta$ is positive, the stratification is unstable for disturbances of all wave numbers exceeding $k_d$ given by
\begin{equation}
\label{eq:10-169}
k_d^2\,d^2 = \frac{4\Omega^2}{g\beta}\,(\tfrac{1}{4}\beta^2d^2+m^2\pi^2),
\end{equation}
while it is stable for disturbances with $k < k_d$. Thus, rotation stabilizes the potentially unstable arrangement for all wave numbers less than
\begin{equation}
\label{eq:10-170}
k_{d,\min}^2 = \frac{4\Omega^2}{g\beta\,d^2}(\tfrac{1}{4}\beta^2d^2+\pi^2).
\end{equation}
We observe that as a function of $\beta$, $k_{d,\min}$ (for a given $\Omega$) attains its minimum value for $\beta = 2\pi/d$: distributions with $\beta$ less than, or greater than, $2\pi/d$ are stabilized by rotation for greater ranges of $k$.

\section{The effect of a vertical magnetic field}\label{sec:96}

We shall now consider the effect of a magnetic field in the direction of $\mathbf{g}$ on the character of the equilibrium of a heterogeneous fluid. We shall suppose that the fluid is inviscid and is of zero resistivity; and we shall not consider the effects of surface tension. It is not difficult to include the effects of finite viscosity, resistivity, and surface tension in the formal analysis; but the essential new elements which a magnetic field introduces can best be appreciated without the complications of many additional parameters which these other effects introduce.

Starting, then, from the equations of hydromagnetics in Chapter IV, \S\,38\,(b), we readily find that under the circumstances postulated the relevant perturbation equations are:
\begin{gather}
\rho\frac{\partial u}{\partial t}-\frac{\mu H}{4\pi}\!\left(\frac{\partial h_x}{\partial z}-\frac{\partial h_z}{\partial x}\right) = -\frac{\partial}{\partial x}\,\delta p, \label{eq:10-171}\\
\rho\frac{\partial v}{\partial t}-\frac{\mu H}{4\pi}\!\left(\frac{\partial h_y}{\partial z}-\frac{\partial h_z}{\partial y}\right) = -\frac{\partial}{\partial y}\,\delta p, \label{eq:10-172}\\
\rho\frac{\partial w}{\partial t} = -\frac{\partial}{\partial z}\,\delta p - g\,\delta\rho, \label{eq:10-173}\\
\frac{\partial\mathbf{h}}{\partial t} = \frac{H}{n}\,D\mathbf{u} \quad\text{and}\quad\frac{\partial}{\partial t}\,\delta\rho = -w\frac{d\rho}{dz}, \label{eq:10-174}
\end{gather}
where $H$ denotes the strength of the uniform magnetic field impressed in the $z$-direction, $\mathbf{h} = (h_x,h_y,h_z)$ is the perturbation in $H$, and the remaining symbols have their usual meanings. In addition to~\eqref{eq:10-171}--\eqref{eq:10-174} we have the equations expressing the solenoidal character of $\mathbf{u}$ and $\mathbf{h}$.

For solutions having the $(x,y,t)$-dependence given by~\eqref{eq:10-13}, the perturbation equations become
\begin{gather}
\rho nu - \frac{\mu H}{4\pi}(Dh_x-ik_xh_z) = -ik_x\,\delta p, \label{eq:10-175}\\
\rho nv - \frac{\mu H}{4\pi}(Dh_y-ik_yh_z) = -ik_y\,\delta p, \label{eq:10-176}\\
\rho nw = -D\,\delta p + \frac{g}{n}(D\rho)w, \label{eq:10-177}\\
h_z = \frac{H}{n}\,Du, \label{eq:10-178}\\
ik_xu+ik_yv = -Dw, \quad\text{and}\quad ik_xh_x+ik_yh_y = -Dh_z. \label{eq:10-179}
\end{gather}
Multiplying equations~\eqref{eq:10-175} and~\eqref{eq:10-176} by $-ik_x$ and $-ik_y$, respectively, adding, and making use of equations~\eqref{eq:10-179}, we obtain
\begin{equation}
\label{eq:10-180}
\rho nDw - \frac{\mu H^2}{4\pi n}(D^2-k^2)h_z = -k^2\delta p;
\end{equation}
but according to~\eqref{eq:10-178}
\begin{equation}
\label{eq:10-181}
h_z = \frac{H}{n}\,Dw.
\end{equation}
Hence
\begin{equation}
\label{eq:10-182}
\rho Dw - \frac{\mu H^2}{4\pi n^2}(D^2-k^2)Dw = -\frac{k^2}{n}\,\delta p.
\end{equation}
Now eliminating $\delta p$ between equations~\eqref{eq:10-177} and~\eqref{eq:10-182}, we obtain
\begin{equation}
\label{eq:10-183}
D(\rho Dw) - \frac{\mu H^2}{4\pi n^2}(D^2-k^2)D^2w = k^2\rho w - \frac{gk^2}{n^2}(D\rho)w.
\end{equation}

If the fluid is confined between two rigid planes at $z = z_1$ and $z_2$ (say), then the boundary conditions which have to be met at these planes are (see Chapter IV, \S\,42\,(a))
\begin{equation}
\label{eq:10-184}
w = 0 \quad\text{and either}\; Dw\;\text{or}\; D^2w = 0\;\text{depending on whether the confining walls are of perfectly conducting or of insulating material.}
\end{equation}

Equation~\eqref{eq:10-183} together with the boundary conditions~\eqref{eq:10-184} leads to the integral formula
\begin{equation}
\label{eq:10-185}
\int\rho\{|Dw|^2+k^2|w|^2\}\,dz + \frac{\mu H^2}{4\pi n^2}\int\{|D^2w|^2+k^2|Dw|^2\}\,dz = \frac{gk^2}{n^2}\int(D\rho)|w|^2\,dz.
\end{equation}

It can be readily shown that this formula provides the basis for a variational formulation of the problem; however, the principal conclusion which we wish to draw from~\eqref{eq:10-185} is that \emph{the characteristic values of $n^2$ are necessarily real}.

If there should be discontinuities in $\rho$, as there would be if fluids of different densities are superposed one over the other, then the conditions which must be met at the interface are
\begin{equation}
\label{eq:10-186}
w\;\text{and}\;\mathbf{h}\;\text{are both continuous.}
\end{equation}
Normally, in a fluid of infinite electrical conductivity, the only condition on $\mathbf{h}$ we can impose is that the normal component of $\mathbf{h}$ ($h_z$ in the present connexion) is continuous. However, since in the present instance the main field is normal to the boundary, the continuity of the \emph{tangential} stresses requires that $h_x$ and $h_y$ are also continuous. The continuity of $h_z$ requires, according to equation~\eqref{eq:10-181}, the continuity of $Dw$; while the continuity of $h_x$ and $h_y$, according to equations~\eqref{eq:10-179}, requires the continuity of $Dh_z$, and, therefore, of $D^2w$. Thus, altogether
\begin{equation}
\label{eq:10-187}
\textit{$w$, $Dw$, and $D^2w$ must be continuous at an interface between two fluids.}
\end{equation}
A further condition follows from equation~\eqref{eq:10-183} by integrating it across the interface. We obtain
\begin{equation}
\label{eq:10-188}
\Delta_s\!\left\{\rho Dw - \frac{\mu H^2}{4\pi n^2}(D^2-k^2)Dw\right\} = -\frac{gk^2}{n^2}\,\Delta_s(\rho)w_s.
\end{equation}

\subsection*{(a) \textit{Two uniform fluids separated by a horizontal boundary: the unstable case}}

Consider the case of two uniform fluids of densities $\rho_1$ and $\rho_2$ separated by a horizontal boundary at $z = 0$. In each of the two regions of constant $\rho$, equation~\eqref{eq:10-185} reduces to
\begin{equation}
\label{eq:10-189}
(D^2-k^2)w - \frac{\mu H^2}{4\pi\rho n^2}(D^2-k^2)D^2w = 0,
\end{equation}
where, for the present, we are suppressing the subscripts distinguishing the two fluids.

The general solution of equation~\eqref{eq:10-189} is a linear combination of the integrals
\begin{equation}
\label{eq:10-190}
e^{\pm kz} \quad\text{and}\quad e^{\pm qz},
\end{equation}
where
\begin{equation}
\label{eq:10-191}
q^2 = \frac{4\pi\rho n^2}{\mu H^2}.
\end{equation}

Now a result of general validity which we have derived is that $n^2$ must be real. Therefore, in~\eqref{eq:10-191} $n^2$ is either positive or negative. This means that the exponent $q$ in the second of the two integrals in~\eqref{eq:10-190} is either real or purely imaginary. In the former case, we can choose a sign for $q$ such that the solution vanishes at $\pm\infty$; in the latter case no such choice is possible. It is therefore clear that the stable ($n^2 < 0$) and the unstable ($n^2 > 0$) cases require separate treatments. Accordingly, we shall first carry through the analysis on the supposition that $n^2$ is positive. This will enable us to distinguish the stable from the unstable stratifications. The treatment (on the supposition $n^2 > 0$) will apply only to the unstable cases; it will not apply to the stable cases which will require a separate treatment.

We shall suppose, then, that
\begin{equation}
\label{eq:10-192}
q = n\sqrt{\frac{4\pi\rho}{\mu H^2}} \quad\text{and}\quad n > 0.
\end{equation}
Returning to the solution of equation~\eqref{eq:10-189}, we first observe that since $w$ must be bounded both when $z \to +\infty$ (in the upper fluid) and $z \to -\infty$ (in the lower fluid), we can write
\begin{equation}
\label{eq:10-193}
\left.
\begin{aligned}
w_1 &= A_1\,e^{+kz}+B_1\,e^{+q_1 z} \quad (z < 0) \\
w_2 &= A_2\,e^{-kz}+B_2\,e^{-q_2 z} \quad (z > 0)
\end{aligned}
\right\}\!,
\end{equation}
where $A_1$, $B_1$, $A_2$, and $B_2$ are constants of integration,
\begin{equation}
\label{eq:10-194}
q_1 = n\sqrt{\frac{4\pi\rho_1}{\mu H^2}}, \quad\text{and}\quad q_2 = n\sqrt{\frac{4\pi\rho_2}{\mu H^2}}.
\end{equation}
In accordance with~\eqref{eq:10-192} we suppose $n > 0$; only then is the foregoing manner of writing the solutions in the two regions ($z > 0$ and $z < 0$) valid.

On the common boundary ($z = 0$) between the two fluids, $w$, $Dw$, and $D^2w$ must be continuous. Applying these conditions to the solution~\eqref{eq:10-193}, we find:
\begin{gather}
A_1+B_1 = A_2+B_2 \quad (= w_0), \label{eq:10-195}\\
kA_1+q_1\,B_1 = -kA_2-q_2\,B_2 \quad (= (Dw)_0), \label{eq:10-196}\\
k^2A_1+q_1^2\,B_1 = k^2A_2+q_2^2\,B_2 \quad (= (D^2w)_0). \label{eq:10-197}
\end{gather}
The remaining condition~\eqref{eq:10-188} now requires
\begin{align}
\label{eq:10-198}
\rho_2\!\left\{Dw_2 - \frac{1}{q_2^2}(D^2-k^2)Dw_2\right\}_{z=0} - \rho_1\!\left\{Dw_1-\frac{1}{q_1^2}(D^2-k^2)Dw_1\right\}_{z=0} \nonumber\\
= -\frac{gk^2}{n^2}(\rho_2-\rho_1)w_0.
\end{align}

Substituting for $w_1$ and $w_2$ from~\eqref{eq:10-193} and simplifying, we are left with
\begin{equation}
\label{eq:10-199}
-\rho_1\!\left(A_1k+\frac{k^2}{q_1}\,B_1\right) - \rho_2\!\left(A_2k+\frac{k^2}{q_2}\,B_2\right) = -\frac{gk^2}{2n^2}(\rho_2-\rho_1)(A_1+B_1+A_2+B_2).
\end{equation}
Equations~\eqref{eq:10-195}--\eqref{eq:10-197} and~\eqref{eq:10-199} lead to the secular equation
\begin{equation}
\label{eq:10-200}
\Delta(q_1,q_2) = \begin{vmatrix}
1 & 1 & -1 & -1 \\[3pt]
k & q_1 & k & q_2 \\[3pt]
k^2 & q_1^2 & -k^2 & -(q_2^2-k^2) \\[3pt]
\frac{1}{2}R-\alpha_1 & \frac{1}{2}R-\alpha_1k/q_1 & \frac{1}{2}R-\alpha_2 & \frac{1}{2}R-\alpha_2k/q_2
\end{vmatrix} = 0,
\end{equation}
where (cf.\ equations~\eqref{eq:10-108} and~\eqref{eq:10-109})
\begin{equation}
\label{eq:10-201}
\alpha_1 = \frac{\rho_1}{\rho_1+\rho_2}, \quad \alpha_2 = \frac{\rho_2}{\rho_1+\rho_2}, \quad\text{and}\quad R = \frac{gk}{n^2}(\alpha_2-\alpha_1).
\end{equation}
The determinant, $\Delta(q_1,q_2)$, can be reduced to give
\begin{align}
\label{eq:10-202}
\Delta(q_1,q_2) &= \begin{vmatrix}
q_1-k & 2k & q_2-k \\
q_1^2-k^2 & 0 & -(q_2^2-k^2) \\
\alpha_1(q_1-k)/q_1 & R-1 & \alpha_2(q_2-k)/q_2
\end{vmatrix} \nonumber\\
&= (q_1-k)(q_2-k)\begin{vmatrix}
q_1+k & 2k & -(q_2+k) \\
\alpha_1/q_1 & R-1 & \alpha_2/q_2
\end{vmatrix} = 0.
\end{align}
Accordingly, $q_1 = k$ and $q_2 = k$ belong to these two characteristic roots; but these two roots lead to trivial solutions: for, one may verify that the proper functions $w_1$ and $w_2$ belonging to these two roots are identically zero. Therefore, removing the factors $q_1-k$ and $q_2-k$ and expanding the remaining determinant, we find
\begin{equation}
\label{eq:10-203}
(R-1)(q_1+q_2+2k) = 2k\!\left[\frac{\alpha_2}{q_2}(q_1+k)+\frac{\alpha_1}{q_1}(q_2+k)\right]\!.
\end{equation}
Define the Alfv\'en velocity
\begin{equation}
\label{eq:10-204}
V_a^2 = \frac{\mu H^2}{4\pi(\rho_1+\rho_2)}.
\end{equation}
In terms of $V_a$,
\begin{equation}
\label{eq:10-205}
q_1 = \frac{n}{V_a}\sqrt{\alpha_1} \quad\text{and}\quad q_2 = \frac{n}{V_a}\sqrt{\alpha_2}.
\end{equation}
Inserting these values in~\eqref{eq:10-203} and substituting also for $R$ in accordance with its definition, we obtain after some minor rearrangements
\begin{equation}
\label{eq:10-206}
\frac{gk}{n^2}(\alpha_2-\alpha_1)\!\left[\frac{k}{\sqrt{\alpha_1}}+\frac{n}{V_a(\sqrt{\alpha_1}+\sqrt{\alpha_2})}\right] = \frac{n}{V_a}+2(\sqrt{\alpha_1}+\sqrt{\alpha_2})+\frac{2kV_a}{n}.
\end{equation}
(It may be noted here for future reference that, in passing from equation~\eqref{eq:10-203} to~\eqref{eq:10-206}, a factor $k/(\sqrt{\alpha_1}+\sqrt{\alpha_2})$ has been removed.)

By measuring $n$ and $k$ in the units
\begin{equation}
\label{eq:10-207}
(g/V_a)\;\text{sec}^{-1} \quad\text{and}\quad (g/V_a^2)\;\text{cm}^{-1},\;\text{respectively,}
\end{equation}
we can rewrite equation~\eqref{eq:10-206} in the non-dimensional form
\begin{equation}
\label{eq:10-208}
\frac{k}{n}(\alpha_2-\alpha_1)\!\left[\frac{k}{\sqrt{\alpha_1}+\sqrt{\alpha_2}}\right] = \frac{n}{k}+2(\sqrt{\alpha_1}+\sqrt{\alpha_2})+2\frac{k}{n}.
\end{equation}
On further reduction, equation~\eqref{eq:10-208} becomes
\begin{equation}
\label{eq:10-209}
n^2+2k(\sqrt{\alpha_1}+\sqrt{\alpha_2})n^2+k(2k+\alpha_1-\alpha_2)n+2k^3(\sqrt{\alpha_2}-\sqrt{\alpha_1}) = 0.
\end{equation}
If $\alpha_2 > \alpha_1$, the constant term in~\eqref{eq:10-209} is negative. The equation must, therefore, allow at least one positive real root; and it can be readily shown that the two remaining roots are either real and negative, or complex conjugates with negative real parts. Since the analysis has presupposed that $n$ is positive, it is clear that the only admissible root of~\eqref{eq:10-209}, in case $\alpha_2 > \alpha_1$, is the single positive root which it allows. The root is

listed in Table~L as a function of $k$ for some values of $\alpha_2$; the derived dispersion relations are further illustrated in Fig.~115.

\begin{table}[ht]
\centering
\caption{The dispersion relations for Rayleigh--Taylor instability when the impressed magnetic field is parallel to $\mathbf{g}$}
\label{tab:L}
\small
\begin{tabular}{c|cc|cc|cc|cc}
\hline
$k$ & \multicolumn{2}{c|}{$\alpha_2 = 1.0$} & \multicolumn{2}{c|}{$\alpha_2 = 0.8$} & \multicolumn{2}{c|}{$\alpha_2 = 0.6$} & \multicolumn{2}{c}{$\alpha_2 = 0.6$} \\
& $n$ & & $n$ & & $n$ & & $n$ & \\
\hline
0 & 0 & & 0 & & 0 & & 0 & \\
0.02 & 0.05796 & & 0.05979 & & 1.0 & 0.69156 & 0.35893 & \\
0.05 & 0.31583 & & 0.14489 & & & 0.77585 & 0.39635 & \\
0.10 & 0.46974 & & 0.26909 & & & 0.86890 & 0.34437 & \\
0.20 & 0.40090 & & 0.24964 & & 2.0 & 0.63890 & 0.49313 & \\
0.30 & 0.33208 & & 0.19544 & & 3.0 & 0.86096 & 0.49313 & \\
0.50 & 0.60060 & & 0.34335 & & 5.0 & 1.01024 & 0.41205 & \\
0.80 & 0.55432 & & 0.34665 & & & 1.00000 & & \\
\hline
\end{tabular}
\end{table}

\figurePlaceholder{115}{The dispersion relations for Rayleigh--Taylor instability when the impressed magnetic field is parallel to $\mathbf{g}$. The abscissa measures the wave number in the unit $(g/V_a^2)\;\text{cm}^{-1}$ and the ordinate measures the growth rate in the unit $(g/V_a)\;\text{sec}^{-1}$, where $V_a$ is the Alfv\'en speed.}

From equation~\eqref{eq:10-209} we derive the following asymptotic relations:
\begin{equation}
\label{eq:10-210}
\left.
\begin{aligned}
n^2 &\to k(\alpha_2-\alpha_1) \quad &(k \to 0)\\
n &\to (\sqrt{\alpha_2}-\sqrt{\alpha_1}) \quad &(k \to \infty)
\end{aligned}
\right\}\!.
\end{equation}
For $k \to 0$ the relation between $n$ and $k$, therefore, approaches the corresponding hydrodynamic relation (cf.\ equation~\eqref{eq:10-128}): \emph{the instability of the very long waves is unaffected by the presence of the magnetic field}. However, in contrast to the hydrodynamic case, $n$ does \emph{not} increase indefinitely with $k$: it tends to a definite limit.

\subsection*{(b) \textit{Two uniform fluids separated by a horizontal boundary: the stable case}}

If $\alpha_1 < \alpha_2$, the constant term in equation~\eqref{eq:10-209} is positive and the equation must allow at least one real negative root. Indeed, the real parts of all three roots must be negative. This follows from the theorem that a necessary and sufficient condition for the real parts of all three roots of the cubic,
\begin{equation}
\label{eq:10-211}
n^3+Bn^2+Cn+D = 0
\end{equation}
to be negative is that $BC > D > 0$. For the cubic~\eqref{eq:10-209}, $D > 0$ (when $\alpha_1 > \alpha_2$) and
\begin{align}
\label{eq:10-212}
BC-D &= 2k^2(\sqrt{\alpha_1}+\sqrt{\alpha_2})(2k+\alpha_1-\alpha_2)-2k^3(\sqrt{\alpha_1}-\sqrt{\alpha_2}) \nonumber\\
&= \frac{2k^2}{\sqrt{\alpha_1}+\sqrt{\alpha_2}}\{[1+2\sqrt{(\alpha_1\alpha_2)}]k+(\alpha_1-\alpha_2)k(\alpha_1\alpha_2)\} > 0.
\end{align}

Since the characteristic equation~\eqref{eq:10-209} allows no positive root when $\alpha_1 > \alpha_2$, we conclude that these cases must be stable and that $n$ must be imaginary. Under these circumstances, the solution for $w$ cannot be written as in~\eqref{eq:10-193} by discarding one of the two integrals $e^{+qz}$ and $e^{-qz}$. Letting $n = i\sigma$ (to emphasize that we are now dealing with an imaginary $n$) we should, in fact, write
\begin{equation}
\label{eq:10-213}
w_1 = A_1\,e^{+kz}+B_1^{(\mathrm{out}:-)}\,e^{i\sigma m_1 z}+B_1^{(\mathrm{inc}:-)}e^{-i\sigma m_1 z} \quad (z < 0)
\end{equation}
and
\begin{equation}
\label{eq:10-214}
w_2 = A_2\,e^{-kz}+B_2^{(\mathrm{out}:+)}\,e^{-i\sigma m_2 z}+B_2^{(\mathrm{inc}:+)}\,e^{+i\sigma m_2 z} \quad (z > 0),
\end{equation}
where $V_1'$ and $V_2'$ denote the Alfv\'en speeds in the two media.

In equation~\eqref{eq:10-213}, we may regard the terms in $\exp(-i\sigma z/V_1')$ and $\exp(+i\sigma z/V_1')$ as representing incoming and outgoing Alfv\'en waves in the first medium. Similarly, in equation~\eqref{eq:10-214} we may regard the terms in $\exp(+i\sigma z/V_2')$ and $\exp(-i\sigma z/V_2')$ as representing incoming and outgoing Alfv\'en waves in the second medium.

By applying the boundary conditions~\eqref{eq:10-187} and~\eqref{eq:10-188} to the solutions~\eqref{eq:10-213} and~\eqref{eq:10-214}, we will obtain a set of four linear equations for the six constants in the solutions. Consequently, the equations can be solved for arbitrary $k$ and $\sigma$; and, moreover, four of the constants $A_1$, $B_1^{(\mathrm{out}:-)}$, $A_2$, and $B_2^{(\mathrm{out}:+)}$ (say), can be expressed linearly in terms of the remaining two constants, $B_1^{(\mathrm{inc}:-)}$ and $B_2^{(\mathrm{inc}:+)}$ (say). This means that \emph{in this case there is no dispersion relation between $k$ and $\sigma$}; and, moreover, the solution for $w$ involves two complex constants $B_1^{(\mathrm{inc}:-)}$ and $B_2^{(\mathrm{inc}:+)}$. The possibility of being able to specify the two latter constants, arbitrarily, ensures our ability to specify, at will, the amplitudes and the phases of the incoming Alfv\'en waves, from $-\infty$ and $+\infty$, in the two media; and the lack of a dispersion relation, similarly, ensures our ability to specify, at will, the frequency of the incoming Alfv\'en waves. Given, then, the amplitudes, the phases, and the frequency of the incoming Alfv\'en waves and given, also, the wave number of the deformed interface, we must expect a unique state of perturbation to emerge; and this is exactly the content of the theory.

In some ways it is remarkable that, while in the absence of a magnetic field we have a valid dispersion relation in all cases, it ceases to exist in the presence of a vertical magnetic field when the circumstances are such as to lead to stability.

\section{The effect of a horizontal magnetic field}\label{sec:97}

We shall now consider the effect of a magnetic field transverse to the direction of $\mathbf{g}$ on the character of the equilibrium of a heterogeneous inviscid fluid of zero resistivity. Except for the change in the direction of the impressed magnetic field relative to $\mathbf{g}$, the physical circumstances of the problem are the same as in \S\,96.

Let the direction of the impressed, uniform magnetic field be along the $x$-axis. Equations~\eqref{eq:10-175}--\eqref{eq:10-178} are then replaced by
\begin{gather}
\rho nu = -ik_x\,\delta p, \label{eq:10-215}\\
\rho nw - \frac{\mu H}{4\pi}(ik_xh_y-ik_yh_x) = -ik_y\,\delta p, \label{eq:10-216}\\
\rho nw - \frac{\mu H}{4\pi}(ik_xh_z-Dh_x) = -D\,\delta p+\frac{g}{n}(D\rho)w, \label{eq:10-217}\\
\mathbf{h} = \frac{ik_x}{n}H\mathbf{u}. \label{eq:10-218}
\end{gather}

Inserting for the components of $\mathbf{h}$ in equations~\eqref{eq:10-216} and~\eqref{eq:10-217} in accordance with equation~\eqref{eq:10-218}, we obtain
\begin{equation}
\label{eq:10-219}
\rho nw - \frac{ik_x}{n}\frac{\mu H^2}{4\pi}\,\zeta = -ik_y\,\delta p
\end{equation}
and
\begin{equation}
\label{eq:10-220}
\rho nw - \frac{ik_x}{n}\frac{\mu H^2}{4\pi}(ik_xw-Du) = -D\,\delta p + \frac{g}{n}(D\rho)w,
\end{equation}
where
\begin{equation}
\label{eq:10-221}
\zeta = ik_yv - ik_xu
\end{equation}
is the $z$-component of the vorticity. Now multiplying equations~\eqref{eq:10-218} and~\eqref{eq:10-219} by $-ik_x$ and $ik_y$, respectively, and adding, we obtain
\begin{equation}
\label{eq:10-222}
\left(\rho n+\frac{k_x^2}{n}\frac{\mu H^2}{4\pi}\right)\zeta = 0
\end{equation}
from which we conclude that
\begin{equation}
\label{eq:10-223}
\zeta = ik_yv - ik_xu = 0.
\end{equation}
Hence, equation~\eqref{eq:10-219} becomes
\begin{equation}
\label{eq:10-224}
\rho nw = -ik_y\,\delta p.
\end{equation}
Equations~\eqref{eq:10-215} and~\eqref{eq:10-224} can be combined to give (on addition after multiplication by $-ik_x$ and $-ik_y$, respectively)
\begin{equation}
\label{eq:10-225}
\rho nDw = -k^2\delta p.
\end{equation}
From equation~\eqref{eq:10-223} and the equation,
\begin{equation}
\label{eq:10-226}
ik_xu+ik_yv = -Dw,
\end{equation}
expressing continuity, we find
\begin{equation}
\label{eq:10-227}
u = \pm\frac{k}{k^2}Dw.
\end{equation}
Inserting this expression for $u$ in equation~\eqref{eq:10-220}, we obtain
\begin{equation}
\label{eq:10-228}
\rho nw - \frac{\mu H^2 k_x^2}{4\pi k^2 n}(D^2-k^2)w = -D\,\delta p + \frac{g}{n}(D\rho)w.
\end{equation}
From equations~\eqref{eq:10-225} and~\eqref{eq:10-228}, $\delta p$ can be eliminated to give
\begin{equation}
\label{eq:10-229}
D(\rho Dw) + \frac{\mu H^2 k_x^2}{4\pi n^2}(D^2-k^2)D^2w - k^2\rho w = -\frac{gk^2}{n^2}(D\rho)w.
\end{equation}

In case we consider solutions for which $k_x = 0$, equation~\eqref{eq:10-229} reduces to
\begin{equation}
\label{eq:10-230}
D(\rho Dw) - k^2\rho w = -\frac{gk^2}{n^2}(D\rho)w;
\end{equation}
this is clearly the same as the equation which obtains in the absence of any field. Hence for disturbances which are independent of the coordinate along the direction of $\mathbf{H}$ (assumed perpendicular to $\mathbf{g}$), the presence of the magnetic field does not in any way affect the development of the Rayleigh--Taylor instability. However, for disturbances with $k_x \neq 0$, there is no longer this identity between the hydrodynamic and the hydromagnetic problems. To see what effect the additional term in $k_x^2$ in equation~\eqref{eq:10-229} has on the problem, consider once again the case of two uniform fluids separated by a horizontal boundary at $z = 0$. We must require that
\begin{equation}
\label{eq:10-231}
w\;\text{and}\; h_x\;\text{are continuous at}\; z = 0.
\end{equation}
According to equation~\eqref{eq:10-218}, the continuity of $w$ ensures the continuity of $h_x$. Also, we can verify that the continuity of the tangential stresses leads to no additional conditions in the first order. Consequently, apart from the condition
\begin{equation}
\label{eq:10-232}
\textit{$w$ is continuous at $z = 0$,}
\end{equation}
we have only to satisfy the condition
\begin{equation}
\label{eq:10-233}
\Delta_s(\rho Dw) + \frac{\mu H^2 k_x^2}{4\pi n^2}\,\Delta_s(D^2w) = -\frac{gk^2}{n^2}(\rho_2-\rho_1)w_0
\end{equation}
which follows from integrating equation~\eqref{eq:10-229} across the interface. Applying the condition~\eqref{eq:10-233} to the solutions given in~\eqref{eq:10-47} and~\eqref{eq:10-48} (which continue to be valid for this problem), we obtain
\begin{equation}
\label{eq:10-234}
n^2 = gk\!\left[\frac{\rho_2-\rho_1}{\rho_2+\rho_1}-\frac{\mu H^2 k_x^2}{2\pi(\rho_2+\rho_1)gk}\right]\!.
\end{equation}
Comparing this equation with equation~\eqref{eq:10-51}, we observe that the additional term in $k_x^2$ in equation~\eqref{eq:10-230} has, for this problem, an effect which is the same as a surface tension. We may, indeed, define an equivalent surface tension by the formula
\begin{equation}
\label{eq:10-235}
T_{\text{eq}} = \frac{\mu H^2}{2\pi k}\cos^2\!\vartheta,
\end{equation}
where $\vartheta$ is the inclination of the wave vector $(k_x, k_y)$ to the direction of $\mathbf{H}$; this surface tension has its origin, clearly, in the tension $\mu H^2/4\pi$ which exists along the lines of force (cf.\ Chapter IV, \S\,37).

\section{The oscillations of a viscous liquid globe}\label{sec:98}

The problems considered in the preceding sections can naturally be extended to a spherical geometry. In view of possible applications, such extensions have more than a formal interest; but in most instances, the problems become of such complexity and involve so many parameters that `elementary' methods of solution seem impracticable. Nevertheless, it will be useful to consider at least one problem for which an analytical solution can be found and which will illustrate the type of difficulties one must confront in the other problems. For this illustrative purpose, we shall consider the oscillations of a viscous liquid globe; this problem is the analogue, in a sphere, of the problem of the gravity waves considered in \S\,94\,(e).

\subsection*{(a) \textit{The perturbation equations and their solution}}

Quite generally, the equations governing the departures from an equilibrium state are
\begin{gather}
\frac{\partial\mathbf{u}}{\partial t} = -\operatorname{grad}\varpi + \nu\operatorname{curl}^2\mathbf{u}, \label{eq:10-236}\\
\operatorname{div}\mathbf{u} = 0, \quad \nabla^2\delta V = 0, \label{eq:10-237}
\end{gather}
and
\begin{equation}
\label{eq:10-238}
\varpi = -\delta V + \delta p/\rho,
\end{equation}
where $\mathbf{u}$ denotes the velocity and $\delta p$ and $\delta V$ are the changes in the pressure and the internal gravitational potential consequent to the perturbation.

Since we are considering departures from an equilibrium spherical state of an incompressible fluid, a normal mode can be specified uniquely in terms of the shape of the deformed surface of the fluid. Suppose, then, that the deformed surface is described by the equation
\begin{equation}
\label{eq:10-239}
r = R+\epsilon Y_l^m(\vartheta,\varphi),
\end{equation}
where $R$ is the radius of the unperturbed sphere, $\epsilon$ is some function of the time (which we shall presently specify to be of the form $\epsilon_0\,e^{-\sigma t}$), and $Y_l^m$ is a spherical harmonic.

For deformations of the surface described by~\eqref{eq:10-239}, the change in the gravitational potential $\delta V$ can be readily found from the corresponding linearized forms of Poisson's and Laplace's equations and the conditions that the potential and its derivative are continuous on~\eqref{eq:10-239}. We find\footnote{The solutions of Poisson's and Laplace's equations which have to be matched are
\[
V_{\text{internal}} = \tfrac{4}{3}\pi G\rho(\tfrac{3}{2}R^2-\tfrac{1}{2}r^2)+\epsilon B r^l Y_l^m
\]
and
\[
V_{\text{external}} = GM/r + \epsilon A r^{-(l+1)} Y_l^m.
\]
The conditions that $V$ and its radial derivative are continuous on~\eqref{eq:10-239} give
\[
A/R^{l+1} = BR^l, \quad\text{and}\quad lBR^{l-1}+(l+1)A/R^{l+2} = \tfrac{4}{3}\pi G\rho R.
\]
From these equations we find
$A = \frac{3}{2} GM R^{l-1}/(2l+1)$ and $B = \frac{3}{2}GM/(2l+1)R^{l+1}$,
and the result quoted follows.}
\begin{equation}
\label{eq:10-240}
\delta V = -\frac{3}{2l+1}\frac{GM}{R^{l+1}}\,r^l\,Y_l^m.
\end{equation}

To determine $\delta p/\rho$, we first observe that by taking the divergence of equation~\eqref{eq:10-237} and remembering that $\nabla^2\delta V = 0$ we have
\begin{equation}
\label{eq:10-241}
\nabla^2(\delta p/\rho) = 0.
\end{equation}
We may, therefore, write
\begin{equation}
\label{eq:10-242}
\frac{\delta p}{\rho} = \epsilon P_l\frac{r^l}{R^l}\,Y_l^m,
\end{equation}
where $P_l$ is some constant undetermined for the present.

We now express $\mathbf{u}$ as expressed in terms of the toroidal and the poloidal functions described in Chapter VI, \S\,56\,(d). In the terminology of these functions, $\operatorname{grad}\varpi$ is a solenoidal vector which is poloidal; and its defining scalar function is
\begin{equation}
\label{eq:10-244}
\epsilon\Pi_0\,r^{l+1}.
\end{equation}

In view of this purely poloidal character of the inhomogeneous term in equation~\eqref{eq:10-236}, it follows that the velocity field $\mathbf{u}$ must also be purely poloidal and specifically of the type given in equation VI~(50).

We now suppose that the dependence on time of the various quantities is given by
\begin{equation}
\label{eq:10-245}
\epsilon_0\,e^{-\sigma t},
\end{equation}
where $\sigma$ is a constant to be determined. (Note the minus sign in this definition.) Then $\mathbf{u}$ can be expressed in terms of a single scalar function $U(r)$ in the manner
\begin{equation}
\label{eq:10-246}
u_r = \epsilon\,e^{-\sigma t}\frac{l(l+1)}{r}\,U(r)Y_l^m, \quad u_\theta = \epsilon\,e^{-\sigma t}\frac{1}{r}\frac{dU}{dr}\frac{\partial Y_l^m}{\partial\vartheta}
\end{equation}
and
\begin{equation}
u_\varphi = \epsilon\,e^{-\sigma t}\frac{1}{r\sin\vartheta}\frac{dU}{dr}\frac{\partial Y_l^m}{\partial\varphi},
\end{equation}
where we have written $\epsilon_0\,e^{-\sigma t}$ in place of $\epsilon$.

$\operatorname{grad}\varpi$ is similarly defined in terms of the scalar function
\begin{equation}
\label{eq:10-247}
\epsilon\,e^{-\sigma t}\,\epsilon_0\,\Pi_0\,r^{l+1}.
\end{equation}

\subsection*{(i) \textit{The solution for the inviscid case: the Kelvin modes}}

Before we proceed to a general discussion of equation~\eqref{eq:10-236}, we shall consider the case $\nu = 0$. Equation~\eqref{eq:10-236} then becomes
\begin{equation}
\label{eq:10-248}
\frac{\partial\mathbf{u}}{\partial t} = -\operatorname{grad}\varpi.
\end{equation}
Replacing each side of this equation by its defining scalar function, we have
\begin{equation}
\label{eq:10-249}
\sigma U(r) = \epsilon_0\,\Pi_0\,r^{l+1},
\end{equation}
so that in this case the radial component of the velocity is given by
\begin{equation}
\label{eq:10-250}
u_r = \frac{\epsilon_0}{\sigma}\Pi_0\,l(l+1)r^{l-1}\,Y_l^m\,e^{-\sigma t}.
\end{equation}

The boundary conditions which must be satisfied in this inviscid case are: (1) the radial component of the velocity given by~\eqref{eq:10-250} must be compatible with the assumed form of the deformed boundary, namely,~\eqref{eq:10-251}; and (2) the pressure must vanish on the same boundary. The first of these conditions clearly requires
\begin{equation}
\label{eq:10-251}
r = R + \epsilon_0\,e^{-\sigma t}Y_l^m;
\end{equation}
\begin{equation}
\label{eq:10-252}
\frac{\epsilon_0}{\sigma}\Pi_0\,l(l+1)R^{l-1} = -\sigma\epsilon_0
\end{equation}
or
\begin{equation}
\label{eq:10-253}
\sigma^2 = -l(l+1)\Pi_0\,R^{l-1}.
\end{equation}
Now the pressure $p_0$ in the unperturbed configuration is given by
\begin{equation}
\label{eq:10-254}
p_0 = \tfrac{2}{3}\pi G\rho^2(R^2-r^2).
\end{equation}
The vanishing of the pressure $p_0+\delta p$ on the perturbed boundary~\eqref{eq:10-251} requires
\begin{equation}
\label{eq:10-255}
-(\tfrac{4}{3}\pi G\rho R)\epsilon Y_l^m + \left(\frac{\delta p}{\rho}\right)_{r=R} = 0.
\end{equation}
On the other hand, according to equation~\eqref{eq:10-243},
\begin{equation}
\label{eq:10-256}
-\frac{3}{2l+1}\frac{GM}{R^2}Y_l^m = \epsilon(l+1)\Pi_0\,R^{l}\,\epsilon Y_l^m.
\end{equation}
Eliminating $\left(\frac{\delta p}{\rho}\right)_{r=R}$ between these equations, we have
\begin{equation}
\label{eq:10-257}
(l+1)\Pi_0\,R^l = \frac{2(l-1)}{2l+1}\frac{GM}{R^2}.
\end{equation}
Finally, combining equations~\eqref{eq:10-253} and~\eqref{eq:10-257}, we obtain
\begin{equation}
\label{eq:10-258}
\sigma^2 = -\frac{2l(l-1)}{2l+1}\frac{GM}{R^3} = \frac{2l(l-1)}{2l+1}\frac{4\pi G\rho}{3}\cdot\frac{l(l-1)}{2l+1}.
\end{equation}
This formula giving the frequencies of oscillation of an inviscid liquid globe is due to Kelvin.

\subsection*{(ii) \textit{The solution for the general case}}

Returning to equation~\eqref{eq:10-236} and replacing each term by its defining scalar function, we have
\begin{equation}
\label{eq:10-259}
-\sigma U = -\epsilon_0\,\Pi_0\,r^{l+1}+\nu\!\left\{\frac{d^2U}{dr^2}-\frac{l(l+1)}{r^2}\,U\right\}\!,
\end{equation}
where we have made use of the result quoted in \S\,56\,(d) that the defining scalar of $\operatorname{curl}^2\mathbf{S}$ is $\bar{S}$ (defined in equation VI~(56)). We rewrite equation~\eqref{eq:10-259} in the form
\begin{equation}
\label{eq:10-260}
\frac{d^2U}{dr^2}-\frac{l(l+1)}{r^2}\,U+\frac{\sigma}{\nu}\,U = \frac{\epsilon_0}{\nu}\,\Pi_0\,r^{l+1}.
\end{equation}
The general solution of this equation which is free of singularity at the origin is
\begin{equation}
\label{eq:10-261}
U = Ar J_{l+\frac{1}{2}}(qr)+\frac{\epsilon_0}{\nu}\,\Pi_0\,r^{l+1},
\end{equation}
where $A$ is a constant, $J_{l+\frac{1}{2}}(z)$ is the Bessel function of order $l+\frac{1}{2}$, and
\begin{equation}
\label{eq:10-262}
q = \sqrt{(\sigma/\nu)}.
\end{equation}
According to equations~\eqref{eq:10-246} and~\eqref{eq:10-261}, the radial component of the velocity is given by
\begin{equation}
\label{eq:10-263}
u_r = l(l+1)\!\left\{A\frac{J_{l+\frac{1}{2}}(qr)}{r^{1/2}}+\frac{\epsilon_0}{\nu}\,\Pi_0\,r^{l-1}\right\}\!Y_l^m\,e^{-\sigma t}.
\end{equation}

\subsection*{(b) \textit{The boundary conditions and the characteristic equation}}

The solution represented by equation~\eqref{eq:10-261} must satisfy certain boundary conditions. These are: (1) the radial component of the velocity given by~\eqref{eq:10-263} must be compatible with the assumed form of the deformed boundary, namely,~\eqref{eq:10-251}; (2) the tangential viscous stresses must vanish at $r = R$; and (3) the $(r,r)$-component of the total stress tensor must vanish on the deformed boundary of the configuration.

The first of the foregoing conditions clearly requires (cf.\ equation~\eqref{eq:10-252})
\begin{equation}
\label{eq:10-264}
l(l+1)\!\left\{A\frac{J_{l+\frac{1}{2}}(z)}{R^{1/2}}+\frac{\epsilon_0}{\nu}\,\Pi_0\,R^{l-1}\right\} = -\epsilon_0\,\sigma,
\end{equation}
where
\begin{equation}
\label{eq:10-265}
z = qR = \sqrt{(\sigma R^2/\nu)}.
\end{equation}

The general expressions for the tangential viscous stresses in spherical polar coordinates have been given in Chapter VI (equation~(44)). The boundary conditions require that these stresses vanish at $r = R$. Substituting for the components of $\mathbf{u}$ from~\eqref{eq:10-246} in equations VI~(44), we verify that the conditions require, only,
\begin{equation}
\label{eq:10-266}
\frac{l(l+1)}{r}\,U - \frac{2}{r}\frac{dU}{dr}+\frac{d^2U}{dr^2} = 0 \quad\text{for}\;r = R.
\end{equation}
Applying this condition to $U$ given by~\eqref{eq:10-261}, we find after some reductions (in which use is made of the various recurrence relations satisfied by the Bessel functions)
\begin{equation}
\label{eq:10-267}
\frac{A}{R^{1/2}}\!\left[2zJ_{l+\frac{1}{2}}(z)-z^2J'_{l+\frac{1}{2}}(z)\right]+2(l^2-1)\!\left\{A\frac{J_{l+\frac{1}{2}}(z)}{R^{1/2}}+\frac{\epsilon_0}{\nu}\,\Pi_0\,R^{l-1}\right\} = 0.
\end{equation}

From equations~\eqref{eq:10-264} and~\eqref{eq:10-267} it follows that
\begin{equation}
\label{eq:10-268}
A = \frac{2(l-1)\epsilon_0\sigma R^{l+1/2}}{l[2zJ_{l+\frac{1}{2}}(z)-z^2J'_{l+\frac{1}{2}}(z)]}.
\end{equation}
where we have introduced the function
\begin{equation}
\label{eq:10-270}
Q_l(z) = J_{l+\frac{1}{2}}(z)/J'_{l+\frac{1}{2}}(z).
\end{equation}

It remains to satisfy the last of the boundary conditions which requires the $(r,r)$-component of the total stress tensor to vanish on the deformed boundary of the configuration. The required component of the stress tensor is given by
\begin{equation}
\label{eq:10-271}
P_{rr} = p_0+\delta p-2\nu\rho\frac{\partial u_r}{\partial r},
\end{equation}
where $p_0$ is the pressure in the unperturbed configuration given by~\eqref{eq:10-254}. Retaining only terms of order $\epsilon$, we find that the vanishing of $P_{rr}$ on $R+\epsilon Y_l^m$ requires
\begin{equation}
\label{eq:10-272}
-(\tfrac{4}{3}\pi G\rho R)\epsilon Y_l^m+\left(\frac{\delta p}{\rho}\right)_{r=R}-2\nu\!\left(\frac{\partial u_r}{\partial r}\right)_{r=R} = 0.
\end{equation}
On the other hand, according to equation~\eqref{eq:10-256}),
\begin{equation}
\label{eq:10-273}
\left\{-\frac{2(l-1)}{2l+1}\frac{GM}{R^2}+\frac{3}{2l+1}\frac{GM}{R^2}+(l+1)\Pi_0\,R^l\right\}\epsilon Y_l^m = 0.
\end{equation}
Inserting this latter expression for $(\delta p/\rho)_{r=R}$ in equation~\eqref{eq:10-272}, we find
\begin{equation}
\label{eq:10-274}
\left\{-\frac{2(l-1)}{2l+1}\frac{GM}{R^2}+l(l+1)\Pi_0\,R^l\right\}\epsilon Y_l^m - 2\nu\!\left(\frac{\partial u_r}{\partial r}\right)_{r=R} = 0.
\end{equation}
This equation can be written more conveniently in the form
\begin{equation}
\label{eq:10-275}
\{-\sigma_{Kl}^2+l(l+1)\Pi_0\,R^{l-1}\}\epsilon Y_l^m - \frac{2\nu}{R}\!\left(\frac{\partial u_r}{\partial r}\right)_{r=R} = \sigma_{Kl}^2,
\end{equation}
where $\sigma_{Kl}$ is the Kelvin frequency (cf.\ equation~\eqref{eq:10-258}).

Now evaluating $\partial u_r/\partial r$ in accordance with equations~\eqref{eq:10-263} and~\eqref{eq:10-268}, we find
\begin{equation}
\label{eq:10-276}
\left(\frac{\partial u_r}{\partial r}\right)_{r=R} = l(l^2-1)\!\left\{\frac{2(l-1)}{lR}\frac{zQ_{l+\frac{1}{2}}(z)}{2zQ_{l+\frac{1}{2}}(z)-z^2}+\Pi_0\,R^{l-1}\right\}\!\epsilon\,e^{-\sigma t}Y_l^m.
\end{equation}
Using this result in equation~\eqref{eq:10-275}, we have
\begin{align}
\label{eq:10-277}
l(l+1)\Pi_0\,R^{l-1}\!\left[1-\frac{2l(l-1)}{z^2}\right] + \frac{4\nu\sigma}{R^2}\frac{l(l^2-1)}{z^2}\cdot\frac{(l-1)-zQ_{l+\frac{1}{2}}(z)}{z^2-2zQ_{l+\frac{1}{2}}(z)} = \sigma_{Kl}^2.
\end{align}

Remembering that $z = \sqrt{(\sigma R^2/\nu)}$, we can rewrite the foregoing equation in the form
\begin{equation}
\label{eq:10-278}
l(l+1)\Pi_0\,R^{l-1}\!\left[1-\frac{2l(l-1)}{z^2}\right]+4\sigma^2\!\left[\frac{l(l^2-1)}{z^4}\cdot\frac{(l-1)-zQ_{l+\frac{1}{2}}(z)}{z^2-2zQ_{l+\frac{1}{2}}(z)}\right] = \sigma_{Kl}^2.
\end{equation}

Now eliminating $\Pi_0$ between equations~\eqref{eq:10-269} and~\eqref{eq:10-278}, we obtain
\begin{equation}
\label{eq:10-279}
\sigma^2\!\left[1-\frac{2l(l-1)}{z^2}\right]\!\left[1+4\sigma^2\!\left(\frac{(l+1)Q_{l+\frac{1}{2}}(z)}{z^2-2zQ_{l+\frac{1}{2}}(z)}\right)\right] = \sigma_{Kl}^2.
\end{equation}
After some simplifications, equation~\eqref{eq:10-279} can be brought to the form
\begin{equation}
\label{eq:10-280}
\frac{\sigma}{\sigma_{Kl}^2} = \frac{2(l^2-1)}{z^2}-1+\frac{2l(l-1)}{z^4}\!\left[1-\frac{(l+1)Q_{l+\frac{1}{2}}(z)}{z^2-Q_{l+\frac{1}{2}}(z)}\right]\!.
\end{equation}
This is the required characteristic equation for $\sigma$; it will determine how viscosity damps Kelvin modes.

\subsection*{(c) \textit{The manner of decay of the Kelvin modes. Numerical results}}

Let
\begin{equation}
\label{eq:10-281}
\sigma^2 = \sigma_{Kl}\,R^2/\nu.
\end{equation}
The parameter $\sigma^2$ is determined by the constitution of the globe; its density $\rho$ (which determines $\sigma_{Kl}$), its viscosity $\nu$, and its radius $R$. Indeed, $\sigma^{-2}$ is a measure of the kinematic viscosity in the unit $(R^2\sigma_{Kl})^{-1}$.

In terms of $\sigma_1$,
\begin{equation}
\label{eq:10-282}
x = \sqrt{(\sigma R^2/\nu)} = \sigma_1\sqrt{(\sigma/\sigma_{Kl})};
\end{equation}
and we may rewrite equation~\eqref{eq:10-280} in the form
\begin{align}
\label{eq:10-283}
\sigma^4 &= x^4\!\left[\frac{2(l^2-1)}{x^2}-1\right]+2l(l-1)x^2\!\left[1+\frac{(l+1)Q_{l+\frac{1}{2}}(x)}{x^2-Q_{l+\frac{1}{2}}(x)}\right] \nonumber\\
&= 2(l-1)x^2\!\left\{(l+1)\frac{x^2-2xQ_{l+\frac{1}{2}}(x)}{x^2-Q_{l+\frac{1}{2}}(x)}\right\} \nonumber\\
&\equiv \Psi(x)\;\text{(say).}
\end{align}
Considering first the aperiodic modes of decay (for which $x$ must be real), we observe that when $x \to 0$, the function $\Psi(x)$ is positive and has the behaviour
\begin{equation}
\label{eq:10-284}
\Psi(x) = 2\,(l-1)(2l^2+4l+3)\,x^2+O(x^4) \quad (x \to 0).
\end{equation}
Hence
\begin{equation}
\label{eq:10-285}
\sigma^4 \sim \frac{2l(l-1)(2l^2+4l+3)}{2l+1}x^2,
\end{equation}
and
\begin{equation}
\label{eq:10-286}
\sigma \sim \sigma^2\frac{2l+1}{2l(l-1)(2l^2+4l+3)} \quad (x \to 0).
\end{equation}
Substituting for $\sigma^2$ from equation~\eqref{eq:10-281} in~\eqref{eq:10-286}, we obtain
\begin{equation}
\label{eq:10-287}
\sigma \to \sigma_{Kl}^2\frac{2l+1}{2(l-1)(2l^2+4l+3)}\cdot\frac{R^2}{\nu} \quad\text{as}\;\nu \to \infty.
\end{equation}

The infinitely slow damping which the foregoing relation predicts for $\nu \to \infty$ means only that in this problem we have the analogues of the `creeping' modes which we have encountered in \S\,94\,(e).

As $x$ increases from zero, $\Psi(x)$ starts increasing, attains a maximum (at which point $\sigma^2$ also attains a maximum $\sigma_{\max}^2$, say), and then decreases to zero at a finite $x$. This is the first of a succession of zeros and there are in reality an infinity of intervals in which $\Psi(x)$ is positive. Postponing the consideration of these other intervals and restricting ourselves for the present to the first interval (including the origin) in which $\Psi(x)$ is positive, we observe that for $\sigma^2 < \sigma_{\max}^2$ there are two real roots of equation~\eqref{eq:10-283} which lead to two aperiodic modes of decay. In this respect, also, the solution for the present problem resembles the solution for the plane problem. These facts are apparent from the numerical results given in Table~LI and illustrated in Fig.~116.

\begin{table}[ht]
\centering
\caption{The aperiodic modes of decay in the oscillations of a viscous liquid globe}
\label{tab:LI}
\small
\begin{tabular}{c|cc|cc|cc}
\hline
& \multicolumn{2}{c|}{$l = 2$} & \multicolumn{2}{c|}{$l = 3$} & \multicolumn{2}{c}{$l = 4$} \\
$\sigma_1^2 R/\nu$ & $\sigma/\sigma_{Kl}$ & & $\sigma/\sigma_{Kl}$ & & $\sigma/\sigma_{Kl}$ & \\
\hline
0 & 0 & & 0 & & 0 & \\
0.2755 & 0.09530 & & 0.09526 & & 0.05970 & 2.1197 \\
1.8630 & 0.26301 & & 0.21809 & & & \\
2.5381 & 0.30669 & & 0.24041 & & & \\
3.9588 & 0.37088 & & 0.30060 & & 0.10288 & \\
3.5927 & 0.39513 & & & & & \\
3.9588 & 0.44600 & & 0.31600 & & & \\
3.5988 & 0.37088 & & & & & \\
3.9685 & 0.43888 & & 0.34657 & & & \\
\hline
\end{tabular}
\end{table}

\figurePlaceholder{116}{The dependence of the damping constant $\sigma$ for the aperiodic modes on the constitution of the globe $R$, $\rho$, and $\nu$. The ordinate measures $\sigma$ in units of the circular frequency $\sigma_{Kl}$ of the undamped $l = 2$ Kelvin mode. The abscissa measures the reciprocal of the kinematic viscosity in the unit $(2R_c\sigma_{Kl})^{-1}$. The curves are labelled by the values of $l$ to which they refer.}

positive, we observe that for $\sigma^2 < \sigma_{\max}^2$ there are two aperiodic modes of decay. For $\sigma^2 > \sigma_{\max}^2$, the lowest modes of decay must be characterized by complex $\sigma$'s with negative real parts. A complete numerical analysis of the corresponding characteristic values is difficult because of the lack of tables of spherical Bessel functions with complex arguments. However, the behaviour of $\sigma$ for $\nu \to 0$ can be readily established: for, as $\nu \to 0$, $|z| \to \infty$ and the asymptotic behaviour of $\sigma$ is determined by (cf.\ equation~\eqref{eq:10-280})
\begin{equation}
\label{eq:10-289}
\frac{\sigma_{Kl}^2}{\sigma^2} = \frac{2l(l-1)(2l+1)}{z^2}-1.
\end{equation}
Combining this relation with the definition of $x$ (equation~\eqref{eq:10-282}), we find
\begin{equation}
\label{eq:10-290}
\sigma = \pm i\sigma_{Kl}+(l-1)(2l+1)\nu/R^2 \quad\text{as}\;\nu \to 0.
\end{equation}
Therefore, in the limit of zero viscosity, the Kelvin modes are damped with mean lives given by
\begin{equation}
\label{eq:10-291}
\tau_l = \frac{R^2}{(l-1)(2l+1)\nu} \quad (\nu \to 0).
\end{equation}

In addition to the lowest modes of aperiodic decay which we have described and illustrated, there are an infinity of other (higher order) modes (with larger damping constants) which can be derived from equation~\eqref{eq:10-283} by considering intervals, beyond the first, in which $\Psi(x)$ is positive. The existence of these other intervals is associated with the poles of $Q_{l+\frac{1}{2}}(x)$ which occur at the zeros, $x_{l+\frac{1}{2}}$, of $J'_{l+\frac{1}{2}}(x)$. In the neighbourhood of each of these poles of $Q_{l+\frac{1}{2}}(x)$ the expression $x^2-2xQ_{l+\frac{1}{2}}(x)$ has a zero; and at these zeros $\Psi(x)$ has poles. To a good approximation, these higher modes of decay are, therefore, characterized by the damping constants
\begin{equation}
\label{eq:10-288}
\sigma \sim x_{l+\frac{1}{2},j}^2\nu/R^2 \quad (j \geqslant 2).
\end{equation}
We have seen that for $\sigma^2 < \sigma_{\max}^2$, there are two principal modes of aperiodic decay. For $\sigma^2 > \sigma_{\max}^2$, the lowest modes of decay must be characterized by complex $\sigma$'s with negative real parts. A complete numerical analysis of the corresponding characteristic values is difficult because of the lack of tables of spherical Bessel functions with complex arguments. However, the behaviour of $\sigma$ for $\nu \to 0$ can be readily established: for, as $\nu \to 0$, $|z| \to \infty$ and the asymptotic behaviour of $\sigma$ is determined by (cf.\ equation~\eqref{eq:10-280}).

\section{The oscillations of a viscous liquid drop}\label{sec:99}

In this section, we shall show that in spite of appearances the solution of the problem of the oscillations of a viscous liquid drop can be reduced to that of the viscous liquid globe considered in \S\,98. The restoring forces maintaining equilibrium in the two cases are very different; in the case of the liquid globe it is gravitation, while in the case of the liquid drop it is surface tension. Nevertheless, it will appear that if we express the time constant $\sigma$ in each case in terms of the appropriate proper frequency $\sigma_{Kl}$, in the absence of viscosity, the characteristic equations determining $\sigma$ become identical. The identity of the two problems was first established by Reid, though in the limit $\nu \to 0$ the theorem was stated by Lamb.

Consider, then, the perturbation of a viscous liquid drop held together by surface tension at the free boundary. Equations~\eqref{eq:10-236} and~\eqref{eq:10-237} governing departures from equilibrium are, of course, of general validity; but in place of equation~\eqref{eq:10-238} we now have
\begin{equation}
\label{eq:10-292}
\varpi = \delta p/\rho,
\end{equation}
since in this case there is no additional contribution derived from a potential.

A normal mode of perturbation of the drop can again be specified by defining the shape of the deformed boundary as by equation~\eqref{eq:10-239}. And it is clear that the arguments leading to the solution~\eqref{eq:10-242} continue to be valid. Accordingly, we can write
\begin{equation}
\label{eq:10-293}
\varpi = \delta p/\rho = \epsilon(l+1)\Pi_0\,r^{l+1}Y_l^m,
\end{equation}
where $\Pi_0$ is a constant. The solution for $\mathbf{u}$ follows exactly as before; in particular, it can be expressed as in equations~\eqref{eq:10-246}).

\subsection*{(a) \textit{The solution for the inviscid case}}

The arguments leading to equation~\eqref{eq:10-253} are unaffected, and we have
\begin{equation}
\label{eq:10-294}
\sigma^2 = -l(l+1)\Pi_0\,R^{l-1}.
\end{equation}
The analysis following equation~\eqref{eq:10-253}, which seeks to determine $\Pi_0$, is not valid now; for, in the unperturbed state, the pressure $p_0$ inside the drop is constant and is given by
\begin{equation}
\label{eq:10-295}
p_0 = 2T/R,
\end{equation}
where $T$ denotes the surface tension. (We are restricting our discussion to the case when the drop is in free space and the external pressure is zero.) According to equation~\eqref{eq:10-24}), the condition which we must satisfy in the perturbed state is that on the deformed boundary,
\begin{equation}
\label{eq:10-296}
p_0+\delta p = T\!\left(\frac{1}{R_1}+\frac{1}{R_2}\right)\!,
\end{equation}
where $R_1$ and $R_2$ are the principal radii of curvature of the surface
\begin{equation}
\label{eq:10-297}
r = R+\epsilon Y_l^m.
\end{equation}
It can be shown\footnote{Cf.\ Lamb, \textit{Hydrodynamics}, Cambridge, England, 1932; p.~475.} that for the surface defined by~\eqref{eq:10-297}
\begin{equation}
\label{eq:10-298}
\frac{1}{R_1}+\frac{1}{R_2} = \frac{2}{R}+\frac{\epsilon}{R^2}(l-1)(l+2)Y_l^m.
\end{equation}
Hence, in this case
\begin{equation}
\label{eq:10-299}
\left(\frac{\delta p}{\rho}\right)_{r=R} = \epsilon(l-1)(l+2)\frac{T}{\rho R^2}\,Y_l^m.
\end{equation}
By comparison with equation~\eqref{eq:10-293} it follows that
\begin{equation}
\label{eq:10-300}
(l+1)\Pi_0\,R^{l-1} = (l-1)(l+2)\frac{T}{\rho R^3}.
\end{equation}
Using this result in equation~\eqref{eq:10-294}, we obtain
\begin{equation}
\label{eq:10-301}
\sigma^2 = -l(l-1)(l+2)\frac{T}{\rho R^3}.
\end{equation}

\subsection*{(b) \textit{The solution for the general case}}

Returning to the general case, we first observe that the arguments leading to equation~\eqref{eq:10-269} continue to be valid; but we must now replace the condition expressed in~\eqref{eq:10-271} by the requirement that
\begin{equation}
\label{eq:10-302}
P_{rr} = p_0+\delta p - 2\nu\rho\frac{\partial u_r}{\partial r} = T\!\left(\frac{1}{R_1}+\frac{1}{R_2}\right)
\end{equation}
on the deformed boundary. By making use of equations~\eqref{eq:10-293},~\eqref{eq:10-295}, and~\eqref{eq:10-298}, equation~\eqref{eq:10-302} gives
\begin{equation}
\label{eq:10-303}
\left\{(l+1)\Pi_0\,R^{l}-(l-1)(l+2)\frac{T}{\rho R^2}\right\}\epsilon Y_l^m - 2\nu\!\left(\frac{\partial u_r}{\partial r}\right)_{r=R} = 0.
\end{equation}
We can rewrite this equation in the form
\begin{equation}
\label{eq:10-304}
\left\{-\sigma_{Kl}^2+l(l+1)\Pi_0\,R^{l-1}\right\}\epsilon Y_l^m - \frac{2\nu}{R}\!\left(\frac{\partial u_r}{\partial r}\right)_{r=R} = 0,
\end{equation}
where (cf.\ equation~\eqref{eq:10-301})
\begin{equation}
\label{eq:10-305}
\sigma_{Kl}^2 = l(l-1)(l+2)\frac{T}{\rho R^3}
\end{equation}
defines the frequencies of oscillation of the normal modes in the absence of viscosity.

Equations~\eqref{eq:10-275} and~\eqref{eq:10-304} are seen to be of identical forms; and since equation~\eqref{eq:10-269} is also true, it is apparent that we shall be led to the same characteristic equation~\eqref{eq:10-280} for $\sigma/\sigma_{Kl}$. The formal identity of the two problems is thus established.

In the examples considered in \S\S\,98 and in the present section, the forces maintaining equilibrium in the unperturbed spherical forms seem to be relevant only for determining the frequencies $\sigma_{Kl}$ of the normal inviscid modes; they do not seem to be significant for determining the manner of their decay by viscosity.

For a drop of water surrounded by air ($T = 74\;\text{dynes/cm}$), this gives $R = 0.23\;\text{mm}$; drops of radii larger than this value will oscillate with decreasing amplitude while drops of smaller radii will be aperiodically damped.

\begin{equation}
\label{eq:10-306}
\sigma_{Kl}\,R^2/\nu = 3.69 \quad\text{and}\quad \sigma^2_{\max}/\sigma_{Kl} = 0.968.
\end{equation}

\section*{BIBLIOGRAPHICAL NOTES}

The following general references may be noted:
\begin{enumerate}
\item \textsc{Sir Horace Lamb}, \textit{Hydrodynamics}, Cambridge, England, 1932.
\item \textsc{J.\ J.\ Proudman}, \textit{Dynamical Oceanography}, John Wiley \& Sons Inc., New York, 1953.
\end{enumerate}

Several of the topics treated in this chapter are considered by Lamb; superposed fluids in \S\,231 (pp.~370--2); the effects of surface tension in \S\S\,266 and 267 (pp.~456--61); gravity waves in \S\,349 (pp.~625--31); and the oscillations of a liquid drop in \S\,275 (pp.~473--5) and \S\,355 (pp.~639--41). Some of these same topics (with interesting oceanographical overtones) are considered by Proudman in chapters XV and XVI of his book (reference~2).

\medskip
\noindent\S\,91. The fundamental paper in the subject is that of Rayleigh:

\begin{enumerate}
\setcounter{enumi}{2}
\item \textsc{Lord Rayleigh}, `Investigation of the character of the equilibrium of an incompressible heavy fluid of variable density', \textit{Scientific Papers}, ii, 200--7, Cambridge, England, 1900.
\end{enumerate}

It is clear that the stability of heterogeneous fluid accelerated in a direction perpendicular to the plane of stratification can be treated by the same formalism: we have merely to replace the acceleration $g$ due to gravity by the actual acceleration $g_a$ (say) to which the fluid is subject. The special case of the stability of the interface between two fluids of differing densities in this latter context was treated by:

\begin{enumerate}
\setcounter{enumi}{3}
\item \textsc{Sir Geoffrey Taylor}, `The instability of liquid surfaces when accelerated in a direction perpendicular to their planes. I', \textit{Proc.\ Roy.\ Soc.\ (London)} A, \textbf{201}, 192--6 (1950).
\end{enumerate}

An experimental demonstration of the development of the `Rayleigh--Taylor' instability (in the case a heavier fluid, overlying a lighter one, is accelerated towards it) is described by:

\begin{enumerate}
\setcounter{enumi}{4}
\item \textsc{D.\ J.\ Lewis}, `The instability of liquid surfaces when accelerated in a direction perpendicular to their planes. II', \textit{Proc.\ Roy.\ Soc.\ (London)} A, \textbf{202}, 81--96 (1950).
\end{enumerate}

The analytical treatment in the text follows largely the presentation in:

\begin{enumerate}
\setcounter{enumi}{5}
\item \textsc{S.\ Chandrasekhar}, `The character of the equilibrium of an incompressible heavy viscous fluid of variable density', \textit{Proc.\ Camb.\ Phil.\ Soc.}\ \textbf{51}, 162--78 (1955).
\end{enumerate}

However, the effects arising from surface tension are not considered in reference~6.

\medskip
\noindent\S\,92. The results described in this section are largely derived from Rayleigh (reference~3). See also the appropriate sections in Lamb (reference~1) and Proudman (reference~2).

\medskip
\noindent\S\,93. The variational formulation given in this section is based on reference~6. Applications of the variational principle to the solution of particular problems will be found in:

\begin{enumerate}
\setcounter{enumi}{6}
\item \textsc{R.\ Hide}, `The character of the equilibrium of an incompressible heavy viscous fluid of variable density: an approximate theory', \textit{Proc.\ Camb.\ Phil.\ Soc.}\ \textbf{51}, 179--201 (1955).
\end{enumerate}

Unfortunately, Hide's particular comparisons with some exact results in reference~6 are vitiated by an oversight to which Reid has drawn attention.

\begin{enumerate}
\setcounter{enumi}{7}
\item \textsc{W.\ H.\ Reid}, `The effects of surface tension and viscosity on the stability of two superposed fluids', \textit{Proc.\ Camb.\ Phil.\ Soc.}\ (in press).
\end{enumerate}

\medskip
\noindent\S\,94. The problem considered in this section was first treated by:

\begin{enumerate}
\setcounter{enumi}{8}
\item \textsc{W.\ J.\ Harrison}, `The influence of viscosity on the oscillations of superposed fluids', \textit{Proc.\ London Math.\ Soc.}\ \textbf{6}, 396--405 (1908).
\end{enumerate}

Harrison's discussion of the problem was quite complete from an analytical standpoint: he derived the basic characteristic equation (equation~\eqref{eq:10-113} in the text); he obtained series expansions for the characteristic values valid in the limit of zero viscosity; and he considered also the effects arising from surface tension. A later discussion of the problem (covering much the same analytical ground) is due to:

\begin{enumerate}
\setcounter{enumi}{9}
\item \textsc{R.\ Bellman} and \textsc{R.\ H.\ Pennington}, `Effects of surface tension and viscosity on Taylor instability', \textit{Quart.\ Appl.\ Math.}\ \textbf{12}, 151--62 (1954).
\end{enumerate}

Detailed numerical results illustrating the effects of viscosity will be found in reference~6. Similarly, the most detailed consideration of the effects arising from surface tension is due to Reid (reference~8). The results quoted and illustrated in this section are taken from these two references.

Discussions on these same and related matters will be found in:

\begin{enumerate}
\setcounter{enumi}{10}
\item \textsc{C.-M.\ Tchen}, `Stability of oscillations of superposed fluids', \textit{J.\ Appl.\ Phys.}\ \textbf{27}, 760--7 (1956).
\end{enumerate}

\begin{enumerate}
\setcounter{enumi}{11}
\item \textsc{G.\ F.\ Carrier} and \textsc{C.\ T.\ Chang}, `On an initial value problem concerning Taylor instability of incompressible fluids', \textit{Quart.\ Appl.\ Math.}\ \textbf{16}, 436--9 (1959).
\end{enumerate}

\begin{enumerate}
\setcounter{enumi}{12}
\item \textsc{S.\ Feldman}, `On the instability theory of the melted surface of an ablating body when entering the atmosphere', \textit{J.\ Fluid Mech.}\ \textbf{6}, 131--55 (1959).
\end{enumerate}

An example of a continuously stratified heterogeneous viscous fluid, in which the density varies exponentially (as in Rayleigh's example considered in \S\,92\,(b)) and the kinematic viscosity remains constant, has been investigated in considerable detail by:

\begin{enumerate}
\setcounter{enumi}{13}
\item \textsc{T.\ Y.\ Teng Fan}, `The character of the instability of an incompressible fluid of constant kinematic viscosity and exponentially varying density', \textit{Astrophys.\ J.}\ \textbf{121}, 508--20 (1955).
\end{enumerate}

A problem which includes features both of heterogeneity and of variation in density due to thermal expansion has recently been considered by:

\begin{enumerate}
\setcounter{enumi}{14}
\item \textsc{H.\ R.\ Morton}, `On the equilibrium of a stratified layer of fluid', \textit{Quart.\ J.\ Mech.\ Appl.\ Math.}\ \textbf{10}, 433--47 (1957).
\end{enumerate}

A related investigation is that of:

\begin{enumerate}
\setcounter{enumi}{15}
\item \textsc{A.\ Skumanich}, `On thermal convection in a polytopic atmosphere', \textit{Astrophys.\ J.}\ \textbf{121}, 408--17 (1955).
\end{enumerate}

\medskip
\noindent\S\,95. The effect of rotation has been considered by:

\begin{enumerate}
\setcounter{enumi}{16}
\item \textsc{R.\ Hide}, `The character of the equilibrium of a heavy, viscous, incompressible, rotating fluid of variable density. I.\ General theory', \textit{Quart.\ J.\ Mech.\ Appl.\ Math.}\ \textbf{9}, 22--34 (1956).

\item[] ---`The character of the equilibrium of a heavy, viscous, incompressible, rotating fluid of variable density. II.\ Two special cases', ibid.\ 35--50 (1956).
\end{enumerate}

Hide's treatment by including the effects of viscosity and of inclination between the directions of $\boldsymbol{\Omega}$ and $\mathbf{g}$ tends to obscure the essential elements of the problem. For this reason, the inviscid case has been treated \textit{de novo}.

\medskip
\noindent\S\,96. The effect of a vertical magnetic field has been considered by:

\begin{enumerate}
\setcounter{enumi}{17}
\item \textsc{R.\ Hide}, `Waves in a heavy, viscous, incompressible, electrically conducting fluid of variable density, in the presence of a magnetic field', \textit{Proc.\ Roy.\ Soc.\ (London)} A, \textbf{233}, 376--96 (1955).
\end{enumerate}

In his treatment of this problem, Hide includes the effects of finite viscosity and resistivity; and the analysis becomes encumbered by many parameters. For the sake of simplicity, the treatment in the text has been restricted to the non-dissipative case.

\medskip
\noindent\S\,97. Kruskal and Schwarzschild first showed that a horizontal magnetic field has no effect on the development of the Rayleigh--Taylor instability:

\begin{enumerate}
\setcounter{enumi}{19}
\item \textsc{M.\ Kruskal} and \textsc{M.\ Schwarzschild}, `Some instabilities of a completely ionized plasma', \textit{Proc.\ Roy.\ Soc.\ (London)} A, \textbf{223}, 348--60 (1954).
\end{enumerate}

\medskip
\noindent\S\,98. The oscillations of a viscous liquid globe were first treated by:

\begin{enumerate}
\setcounter{enumi}{20}
\item \textsc{H.\ Lamb}, `On the oscillations of a viscous spheroid', \textit{Proc.\ London Math.\ Soc.}\ \textbf{13}, 51--66 (1881).
\end{enumerate}

The treatment in the text, however, follows:

\begin{enumerate}
\setcounter{enumi}{21}
\item \textsc{S.\ Chandrasekhar}, `Two oscillations of a viscous liquid globe', ibid.\ (3) \textbf{9}, 141--9 (1959).
\end{enumerate}

The numerical results given in Table~LI are from this reference. Related problems are considered in:

\begin{enumerate}
\setcounter{enumi}{22}
\item \textsc{S.\ Chandrasekhar}, `The character of the equilibrium of an incompressible fluid sphere of variable density and viscosity subject to radial acceleration', \textit{Quart.\ J.\ Mech.\ Appl.\ Math.}\ \textbf{8}, 1--21 (1955).

\item \textsc{W.\ H.\ Reid}, `The oscillations of a viscous liquid globe with a core', \textit{Proc.\ London Math.\ Soc.}\ (3) \textbf{9}, 388--96 (1959).
\end{enumerate}

\medskip
\noindent\S\,99. The principal result established in this section is due to:

\begin{enumerate}
\setcounter{enumi}{24}
\item \textsc{W.\ H.\ Reid}, `The oscillations of a viscous liquid drop', \textit{Quart.\ Appl.\ Math.}\ \textbf{18}, 86--89 (1960).
\end{enumerate}
